\chapter{Results}
\label{ch:analysis}

\section{Robust Black-Box DAG Combiner}
Intuitively, a robust black-box combiner operates on $n$ candidate schemes by 
interleaving oracle calls to the candidates implemented functions with processing 
to combine their outputs in a way such that security properties of the secure 
subset of oracles are preserved if there are sufficiently many of them.
To argue about the impossibility of such constructions it is helpful to impose 
additional structure on the combiner, to structure it as a directed acyclic graph 
modeling the computation of the combiner.
In the DAG every oracle query corresponds to a vertex colored with the appropriate label
Viewing a robust black-box combiner as a robust black-box DAG combiner is therefore 
simply a way make the structured execution of a combiner explicit, making the 
reasoning about oracle query dependencies and oracle-independent computation explicit.
Additionally since robust black-box combiners need to satisfy the security properties relative 
to a PSPACE oracle, in the later proof we rely on the ability of the PSPACE oracle to brute-force 
many secret keys in the algorithm if there are not queries to oracles. For this line of reasoning 
this description is also more convenient. As we will show later in \Cref{thm:generality_dag_comb},
the notions of robust black-box combiners and robust black-box DAG combiners are equivalent.

We define robust black-box DAG combiners by representing each implemented algorithm $\alg{alg}{}$
with a directed acyclic graph $\G_{\alg{alg}{}} = (V_{alg}, E_{alg})$ with a source $s_{\alg{alg}{}} \in V_{\alg{alg}{}}$ 
that has no in-going edges, corresponding to the input of the algorithm, and a sink $t_{alg} \in V_{alg}$ with no out-going edges,
corresponding to a final processing step. Additionally the description of the DAG combiner contains a 
vertex labling $l_{alg}: V_{alg} \setminus \{s_{alg}, t_{alg}\} \to [n] \times \{alg_1, \ldots, alg_k\}$ which
maps every vertex that is not a source or a sink to the corresponding candidate oracle and the specific function 
called on an input $\x_v$ to produce an output $\y_v$. We define the output 
$\y_{s_{alg}}$ of the source to be the combiner algorithm input and the output of the
sink $t_{alg}$ to be the output of the combiner. The input $\x_v$ to the query at vertex $v \in V_{alg} \setminus \{s_{alg}, t_{alg}\}$
is computed by an efficient (ppt.) Turing machine $f_{alg}^v$ associated with vertex $v$ in $G_{alg}$ which takes as input the labeled outputs of all predecessors 
$w \in \mathcal N_{in}(v)$ and produces an oracle query $\x_v$. We formally describe robust black-box DAG combiners 
in the following definition.

\begin{definition}[Robust $(n;k)$ Black-Box DAG $\prim{P}$-Combiner]\label{def:robust_blackbox_dag_combiner}
	Given a cryptographic primitive $\prim{P}$ implementing a set of algorithms
	indexed by $j \in [|\prim{P}|]$. A robust $(n;k)$ black-box DAG $\prim{P}$-combiner 
	is a $|\prim{P}|$-tuple 

	\[
		\Paren*{\mathfrak G_j = \Paren*{\G_j, c_j, s_j, t_j, \{f_{j,v}\}_{v \in V(\G_j)}}}_{j \in [|\prim{P}|]}
	\]

	where $\G_j$ is a directed acyclic graph with vertex set $V(\G_j)$, 
	$c_j: V(\G_j) \setminus \{s_j, t_s\} \to [n] \times [|\prim{P}|]$ is a vertex-coloring with colors
	$(a,b)$ indicating that the corresponding query is made to algorithm $a$ of candidate oracle $b$, source and sink $s_j, t_s \in V(\G_j)$ satisfying $\mathcal N^-(s_j) = \mathcal N^+(t_j) = \varnothing$ and for simplicity
	the only vertices satisfying those properties \todo, and finally $\{f_{j,v}\}_{v \in V(\G_j)}$ are ppt. Turing machines.
	satisfying that the following tuple of oracle Turing machines is a robust $(n;k)$ black-box combiner.

	\[
		\widetilde{\BBComb{\prim{P}}[\O_1, \ldots, \O_n]} = \Paren*{\widetilde{\BBComb{\prim{P}}[\O_1, \ldots, \O_n]}.\interp{\mathfrak G_1}, \ldots, \widetilde{\BBComb{\prim{P}}[\O_1, \ldots, \O_n]}.\interp{\mathfrak G_{|\prim{P}|}}}
	\]
	
	Where $\interp{\mathfrak{G_j}}$ is the interpretation according to \Cref{fig:dag_combiner}.


	\begin{figure}
		\centering 
		\begin{pchstack}[center, space=1em]
			\procedure[linenumbering]
			{Algorithm $\interp{\G, c, s, t, \{f_v\}_{v \in V(\G)}; \rho}(\x)$}
			{
				\y_s \gets \x\\ 
				\pcfor v \in \tau(\G, s) \setminus \{s,t\} \pcdo\\ 
				\t \x_v \gets f_v(\x, \{(\y_w, w)\}_{w \in \mathcal N^-(v)}; \rho)\\ 
				\t \y_v \gets \mathcal O_{c(v)}(\x_v)\\ 
				\pcreturn f_t(\x, \{\y_w\}_{w \in \mathcal N^-(t)})
			}
			\procedure[linenumbering]
			{Algorithm $\tau(\G, s)$}
			{
				\pclinecomment{Topological ordering TODO}
			}
		\end{pchstack}
		\caption{Interpretation of $\Paren*{\G, c, s, t, \{f_v\}_{v \in V(\G)}}$ into an algorithm. $\tau$ returns the vertices in $\G$ in topological order starting from the source $s$.}
		\label{fig:dag_combiner}
	\end{figure}
\end{definition}

A robust black-box DAG combiner algorithm is then executed by tracing the computation through the corresponding DAG $G_{alg}$, 
starting at the source $s_{alg}$ and ending at the sink $t_{alg}$ in a topological ordering. The interpretation of an algorithm 
$alg$ described by the robust black-box DAG combiner by $(\G_{alg}, s_{alg}, t_{alg}, l_{alg}, \{f_{alg}^v\}_{v \in V_{alg} \setminus \{s_{alg}, t_{alg}\}})$
is shown in \Cref{fig:dag_combiner}.

\paragraph{Example.}
We introduce the notion of robust black-box DAG combiners to make the underlying structure of 
robust black-box combiners more explicit. To illustrate the definition we will show the structure of the 
DAG combiner of a robust $(2;1)$ black-box combiner \Cref{def:prf-combiner} in \Cref{fig:example_dag_combiner}.
We label every vertices $v$ in the DAGs that is not a source or sink with $l_{alg}(v)$, i.e. their vertex labling
and for every vertex add the functional description of the Turing machines $f_{alg}^v$.

\begin{figure}[H]
	\centering
	\begin{tikzpicture}
		\node[] (Gen0) at (0,0) {$s_{gen}$};
		\node[] (Gen1) at (-1.25,-1) {$(0, gen), \y_{s_{gen}} \mapsto $};
		\node[] (Gen2) at (1.25,-1) {$(1, gen)$};
		\node[] (Gen3) at (0,-2) {$t_{gen}$};
	\end{tikzpicture}
	\label{fig:example_dag_combiner}
	\caption{Drawing of the graphs $\G$}
\end{figure}

\newpage

Since later on we will focus on showing the impossibility of certain combiner constructions it's 
essential that the notion of combiner that we use is as strong as possible.
By definition a robust black-box DAG combiner is also a robust black-box combiner. 
We show that the notions of robust black-box DAG combiner and robust black-box combiner 
are equally strong and capture the same set of constructions.
We formulate this result in the following theorem.

\begin{theorem}[Generality of Robust $(n;k)$ Black-Box DAG Combiners]\label{thm:generality_dag_comb}
	For every robust $(n;k)$ black-box combiner there exists a computationally equivalent 
	robust $(n;k)$ black-box DAG combiner.   
\end{theorem}

\paragraph{Intuition.}
We need to show that for every robust $(n;k)$ black-box combiner $\BBComb{}[]$ we can construct a 
robuts $(n;k)$ black-box DAG combiner $\widetilde\Pi$ that represents the same function 
if interpreted, given that $\BBComb{}[]$ is efficient is also efficient.
The proof works by unrolling the computation done by the combiner $\BBComb{}[]$ for every oracle 
call made to a candidate scheme.
Since we know that $\BBComb{}[]$ is a black-box combiner, it runs in time polynomial in $\lambda, n, m$. 
Let $T \in \poly[\lambda, n, m]$ be an upper bound on the number of steps that the algorithm runs over all 
inputs $\x$ into the combiner.
We start by constructing a graph $\G_{alg}$ that consists of $T$ layers of $nm$ vertices, 
where $n$ is the number of candidates and $m$ is the number of algorithms implemented by the primitive.
On every layer, each vertex corresponds to an oracle call to a specific oracle algorithm of a candidate.
We connect two vertices $v,w \in V_{alg}$ if the they are not in the same layer with an edge pointing 
from the vertex in the layer closer to the source to the other vertex. To define the functions $f_{alg}^v$
for $v \in V_{alg}$, a vertex in layer $k \in [T]$, we construct a Turing machine that simulates the robust black-box combiner $\BBComb{}[]$ until 
the $k$'th oracle query. Once the $i \in [k-1]$'th oracle query is made we instead use the correct output of the correct vertex in the 
$i$'th layer and move on. For the $k$'th query we simply query the oracle of our vertex if the oracle is correct (and also if its is not correct,
this will then simply be ignored later one during the execution).

To prove that this construction results in the same output as the robust black-box combiner does,
we proceed by induction on the layer depth of the current query, claiming that the output of the 
correct vertex at that layer aligns with the query output of the corresponding query in the execution 
of the robust black-box combiner. Note that the DAG only contains $T \times n \times m$ many vertices and Turing machines $f_{alg}^v$.
Since the computation at every vertex is bounded from above by the time it takes the robust black-box combiner 
to run and we only make one oracle call at every vertex, the result of every vertex can be computed in time at most 
$T$. Overall the time it takes for the robust black-box DAG combiner to run is therefore polynomial in $\lambda, n,m$.
We provide a formal proof below.

\begin{proof}[\Cref{thm:generality_dag_comb}]
	Let $\prim{P}$ be a cryptographic primitive $\BBComb{P}$ with 
	defining algorithms $(p_1, \ldots, p_{|\prim{P}|})$, et $\BBComb{P}$ be a robust $(n;k)$ black-box $\prim{P}$-combiner 
	according to \Cref{def:robust_blackbox_combiner}, and let $j \in [|\prim{P}|]$ be arbitrary.
	By \Cref{def:robust_blackbox_combiner} there exists a polynomial $T \in \poly[n,\lambda,|x|]$ where $x$ is the input to $p_j$.
	We construct a directed acyclic graph $\G_j$ 
	
	\[
		\G_j = \Paren*{[n] \times [|\prim{P}|] \times [T] \cup \{s,t\}, \SetBuilder{\Paren*{(a_1,b_1,t_1), (a_2,b_2,t_2)}}{t_1 < t_2}}
	\]
	
	\todo[and connect to s,t]. We define the coloring by $c(a,b,t) = (a,b)$. We define the Turing machines $\{f_{j,v}\}_{v \in V(\G_j)}$ iteratively starting
	from the beginning of the computation $s$, we define $f_{j,v_{j,a,b,1}}$ to be the ppt Turing machine that simulates the $\BBComb P.p_j$ until the first call 
	to its oracle. If that oracle happens to be $(a,b)$ then we end, otherwise write $0$ onto the tape. In every subsequent level $f_{j,a,b,t+1}$ simulates the combiner 
	on and answers the queries to the respective oracle with the output of that query ($y_{j,a,b,k}$ for $k \in [t]$) by progressively simulating all previous Turing machines 
	$f_{j,a,b,k}$ for $k \in [t]$. By checking whether after the computation we find a $0$ on the tape we can trace the path of queried oracles. We can now continue the simulation of 
	the combiner since we know exactly the state of the tape. \todo[maybe inductive correctness proof but this seems to be very straight forward].
	The combiner is necessarily efficient since for every the number of vertices in the graph is bounded from above by $nT(n,\lambda,|x|)|\prim{P}|$ where 
	$T \in \poly[m,\lambda,|x|]$, we only make exactly one query to any oracle per vertex (which also puts a polynomial bound on the number of oracle calls), 
	the topological ordering can be computed in polynomial time. We are only left to show that the Turing machines $\{f_{j,v}\}_{v \in V(\G_j)}$ are indeed ppt.
	To this end note that $f_{j,v}$ simulates the black-box combiner Turing machine, which is assumed to be ppt, at most once 
	where it can answer any oracle query by looking up the corresponding output of the layer in time polynomial in $n,|\prim{P}|$ and $\lambda$.
\end{proof}

\newpage

\section{Impossibility of Robust Black-Box VRF-Combiner}

We now tackle the main challenge of combining VRFs in a black-box fashion. Specifically 
we study the possibility of constructing a robust $(n;k)$ black-box combiner $\BBComb{VRF}$
for the VRF primitive. During this section we will make the additional assumption that the secret key
space of all candidate schemes is equivalent (or at least is in bijective correspondence to) $\bin^{s(\lambda)}$
for $s$ a polynomial in $\lambda$. \todo{more introduction}

\paragraph{Warm Up.} 
As a warm-up we will start by showing the impossibility of a robust $(2;1)$ black-box combiner for VRFs.
That is, we show that there exists no efficient (ppt.) Turing machine $\BBComb{VRF}$ such that 
for any two oracles $\O_1, \O_2$ implementing a VF, $\BBComb{VRF}[\O_1, \O_2]$ implements a secure VRF 
if at least one of the candidates implements a secure VRF. We state the following lemma.

\begin{lemma}[Robust $(2;1)$ Black-Box VRF Combiner]\label{lemma:warmup}
	There exists no robust $(2;1)$ black-box combiner in the sense of 
	\Cref{def:robust_blackbox_combiner} for the VRF primitive defined in \Cref{def:vrf}.
\end{lemma}

\paragraph{Intuition.}
We must show that for every construction $\BBComb{VRF}$ implemented by a ppt. Turing machine there 
exist two oracles $\O_1, \O_2$, of which at least one implements a secure VRF, such that 
$\BBComb{VRF}$ is either not pseudorandom in the sense of \Cref{def:pseudorandomness},
perfectly correct in the sense of \Cref{def:correctness} or not 
uniquely provable in the sense of  \Cref{def:unique_provability}. 
Applying the result from \Cref{thm:generality_dag_comb} we can argue about the non-existence of DAG combiners 
instead, fix an arbitrary DAG combiner $\BBComb{VRF}$ and consider the DAG corresponding to the evaluation 
algorithm $\G_{eval}$. 
We assume that the combiner is uniquely provable and perfectly correct and construct an adversary
$\A$ that wins the prf security experiment against the combiner with significant advantage.
For that, the adversary will employ the PSPACE oracle $\PSPACE$ to compute a set of candidate secret keys 
by iteratively filtering the secret key space on a polynomial number of evaluation points the adversary 
obtains in the first stage of the prf experiment.
To execute the filtering the adversary constructs a ppt. Turing machine $eval$ that takes as input 
a tuple $(\sk^*, x)$, the combiner input, and returns $(y,\pi)$ such that the output aligns with the combiner output 
if $\sk^*$ is indeed the secret key input and with high probability if the $\sk^*$ is a different secret key.
Importantly $eval$ must not make any queries to the candidate oracles, that way we can use the PSPACE complete oracle 
and an a tuple $(x,y,\pi)$ to compute:

\[
	\SK' = \Brace*{\sk \in \SK : eval^{\PSPACE}(\sk, x) = (y,\pi)}
\]

Assuming that we can efficiently compute $eval$ we can refine the set of candidates $2s(\lambda)$ times.
Now assume that we eliminate at least half the candidates in every second iteration, then we end up with a polynomial set 
of candidates, in which the true secret key is contained. The adversary can now guess the correct secret key 
with non-negligible probability and win the pseudorandomness experiment with non-negligible advantage.
If on the other hand the set of candidates doesnt shrink at half in at least every second step then 
intuitively on random inputs $x$ a lot of secret keys (whether or not they are the correct secret key)
result in the same output of $eval$, which with non-negligible probability agrees with the combiner. 
We can still simply sample a secret key of this set of candidates to obtain a secret key $\sk$,
with possible negligible chance to be the correct secret key, but non-negligible probability to result in the 
same output $y$ on a random $x$. This again enables the adversary to beat the prf security experiment 
with non-negligible advantage.

We now show how to construct the Turing machine $eval$, for that we first fix arbitrary oracles $\O_1, \O_2$ 
which are uniquely provable and perfectly correct, a secret key $\sk$ in the support of the key generation 
function and any evaluation point $x$. We now consider the vertices $v \in V_{eval}$ in a fixed topological 
ordering and classify them for \textit{relevance} in the sense that a vertex is \textit{relevant} if there 
exist two outputs $\hat \y_v \neq \y_v$ in the set of possible outputs of the corresponding oracle query 
(e.g. $\Y_{\O_1} \times \P_{\O_1}$ if the query is to evaluation of candidate $\O_1$)
such that if we simulate the the execution of the combiners evaluation algorithm but replace the 
reply of the oracle with $\hat\y_v$ then we obtain a different output $\hat \y_{t_{eval}} \neq \y_{t_{eval}}$.
We will for now assume that that every oracle query is only queried once on every distinct input $\x$,
we'll show that this is the case later. The same classification of vertices we do for the verification algorithm.

Now initialize the second candidate with $\O_0$ and the first candidate with any secure VRF. To replace the oracle queries in the combiner and arrive at $eval$ on a high level we argue that 
a vertex is relevant in the evaluation algorithm if and only if there exists a relevant vertex in the 
verification algorithm such that the queries made at that vertex is exactly the same as the one done in eval.
Since every relevant query is repeated in the verification algorithm and every not repeated query 
in the verification is irrelevant we can simply define $eval$ to be the Turing machine for $(x,y,\pi,\tau)$
where $\tau$ is the verification algorithms transcript on $x,y,\pi$ that simulates the combiners evaluation algorith,
if he encounters an oracle query to $\O_1$ the simulator checks the transcript $\tau$ if it contains the answer and answers with 
the answer if it exists. If it doesnt exist this means that either $1.$ the secret key doesn't correspond to $(x,y,\pi)$
or and the query is possibly even relevant or $2.$ the secret key is correct, or related, and the query is irrelevant.
If it encounters a query to $\O_2$ it the answer is known.

If we are in case $2.$, then we can simply iterate over the vertices in topological ordering and try every value if the query is not 
in the transcript until every combination of possibile outputs is tested, obtaining the same output every time. 
Assume that were not the case then we could again break unique provability by making $\O_2$ secure again (vrf-ification)
and break unique provability. If we are in case $1.$ then and for every guessing through the combiner (except for the recordings in $\tau$)
the output is $y$, then it is also $y$ for the honest evaluation. We conclude that by only allowing those secret keys to survive at every step,
we conserve \textit{exactly} all secret keys which evaluate $x$ to $y,\pi$.

Now we quickly need to argue that oracle queries are \textit{unique} in the sense that a given algorithm of 
a given candidate is queried at most once on a specific input $\x$ and that we can construct a secure VRF 
from a uniquely provable and perfectly correct VRF that agrees on the original VRF on all secret keys in a polynomial set.
For uniqueness the case is clear since the oracles need to be deterministic, assume they were not deterministic and that two 
outputs in the support would result in a different output, than we would arrive at a contradiction since by perfect correctness both must verify 
since they are honestly generated outputs and by unique provability not both must verify since they are distinct. We can therefore 
assume that the candidates are deterministic or if they are random we can just replace the random output with an arbitrary but fixed output 
in the support. The DAG combiner can now be slightly modified to store the all queries made, if a query is made a second time 
then instead of requerying he could just simply use the result from the first query (since in our DAG construction the 
Turing machine has access to \textit{all} previous outputs).

For the VRF-ification, we give some intuition why a secure VRF can be remapped on a polynomial set of secret keys 
to assume the same output as a uniquely provable and perfectly correct VRF. For that remapping we introduce a symmetric and reflexive relation 
$\sim \subseteq \SK \times \SK$ defined as follows:

\[
	\sk_1 \sim \sk_2 \iff \exists \pk \in \PK, \Prob*{
		\begin{aligned}
			(\sk_1, \pk) &\sample \alg{Gen}{VRF}(1^\lambda) \land \\ (\sk_2, \pk) &\sample \alg{Gen}{VRF}(1^\lambda)
		\end{aligned}
	} > 0
\]

Let $\sim^* := \bigcup_{k \in \N} \sim^k$ be the transitive closure of $\sim$, since $\sim$ is reflexive and symmetric already 
$\sim^*$ is an equivalence relation. We claim that if two keys $\sk_1, \sk_2$ are in the same equivalence relation, i.e. $\sk_1 \sim^* \sk_2$
then $\alg{Eval}{VRF}(\sk_1, \cdot) = \alg{Eval}{VRF}(\sk_1, \cdot)$ since otherwise we could break unique provability. 
If we now want to remap the outputs of a VRF $\O_{\VRF}$ on the set $\widetilde\SK$ then we need to remap the outputs of evaluation in 
$\widehat\SK := \bigcup_{\widetilde\sk \in \widetilde\SK} [\widetilde\sk]_{\sim^*}$, the union of all equivalence classes of $\widetilde\SK$.
Since the newly mapped to points come from a uniquely provable and perfectly correct VF and the secure VRF is also a uniquely provable and 
perfectly correct VRF the resulting VF is uniquely provable and perfectly correct. To show that its pseudorandom we show that in the execution of the pseudorandomness 
game the probability of sampling a secret key $\widehat\sk \in \widehat\SK$ is negligible if $|\widetilde\SK| \in \poly$.
Assume towards a contradiction that the following holds:

\[
	\Prob*{\sk \in \widehat\SK : (\sk, \pk) \sample \alg{Gen}{VRF}(1^\lambda)} > \frac{1}{p(\lambda)}
\]

For a polynomial $p \in \poly$, then we could simply resample a secret key $\sk'$ which would also with probability 
at least $1/p(\lambda)$ be in $\widehat\SK$. Additionally with probability at least $1/|\widehat{\SK}| \in \poly$
we have $\sk \sim^* \sk'$ which means that with probability $1/p(\lambda)|\widehat\SK| \in \poly$
we obtain a secret key which allows us to recompute all outputs, breaking pseudorandomness.

\newpage

Let $\O_1, \O_2$ be arbitrary, but fixed, VF oracles which are uniquely provable 
and perfectly correct and let $\sk, x$ be arbitrary but fixed inputs to the combiners evaluation function, where $\sk$ is in the support 
of the combiners key-generation function \todo{maybe tweak} and $x \in \bin^\lambda$. We iterate over the vertices of $\G_{eval}$ in topological ordering and distinguish two cases,
either the oracle query produced there is irrelevant in the combiners computation of the output in the sense that for 
any output in the support of the oracle (possible output).

If that is not the case there exist possible output $\hat\y_v$ such that after substituting
the output of the oracle with this value the combiner output changes from $\y$ to $\hat\y \neq \y$.
We argue that in this case the exact query on $\x_v$ that is done in the evaluation algorithm 
\textit{must} be repeated in the execution of the verification algorithm. W.l.o.g. assume that 
the query corresponds to oracle $\O_1$. If that were not the case then we 
can construct oracles $\hat\O_1, \O_2^*$ such that $\hat\O_1$ is equivalent to $\O_1$ except that for 
the evaluation in $\x_v$ where it evaluates to $\hat\y_v$ instead of $\y_v$ \todo{first argue that every
oracle query is unique, that is we dont requery the same oracle on the same input, simply build a caching algorithm
and argue using determinism}. We construct $\O_2^*$ to be a secure VRF oracle that agrees with $\O_2$
on all inputs queried during the execution of the combiner on $\sk, x$.
This is possible since we only make a polynomial number of queries to $\O_2$
and even fixing all of the corresponding secret keys (and additional secret keys to preserve unique provability and correctness)
the oracle can still be pesudorandom since the space of possible secret keys is much larger.
Specifically we show that if rewiring those keys leads to an exploit then we could simply guess those keys
and mount the attack anyway. \todo{polish}

Now that we have the oracles $\hat\O_1, \hat\O_2$, of which one is secure, we know that 
$\BBComb{VRF}[\hat\O_1, \hat\O_2]$ and $\BBComb{VRF}[\O_1, \hat\O_2]$ must implement secure VRFs.
In particular by perfect correctness and unique provability we must have the following:
\begin{align}
	\BBComb{VRF}[\hat\O_1, \hat\O_2].\alg{Vfy}{VRF}(\pk, x, y, \pi) &= 1\\
	\BBComb{VRF}[\O_1, \hat\O_2].\alg{Vfy}{VRF}(\pk, x, \hat y, \hat\pi) &= 1\\
	\BBComb{VRF}[\hat\O_1, \hat\O_2].\alg{Vfy}{VRF}(\pk, x, \hat y, \hat \pi) &= 0\\
	\BBComb{VRF}[\O_1, \hat\O_2].\alg{Vfy}{VRF}(\pk, x, y, \pi) &= 0
\end{align}
Since however the oracles $\hat\O_1, \O_1$ only differ in a single evaluation query 
which is not queried by the verification algorithm the evaluation of $\BBComb{VRF}[\hat\O_1, \hat\O_2].\alg{Vfy}{VRF}(\pk, x, y, \pi)$
must evaluate to the same as the evaluation of $\BBComb{VRF}[\hat\O_1, \hat\O_2].\alg{Vfy}{VRF}(\pk, x, \hat y, \hat\pi)$
which is a contradiction. We conclude that the combiner is either not uniquely provable or not 
perfectly corect if there exist irrelevant queries in the evaluation algorithm which are not requeried.
We can additionally show that the converse holds with an analogous argument, establishing a bijection between 
relevant queries in both the evaluation and verification algorithm. 

Directing our attention to the irrelvant queries. As previously argued those queries are \textit{individually irrelevant},
they are also \textit{jointly irrelevant}. We an now instantiate the combiner with two arbitrary \todo{study in which aspect we couldnt just start with VRFs but start wiht uniquely provable and perfectly correct oracles.}
oracles $\O_0, \O_{VRF}$ and argue that given a tuple $(x, y, \pi)$ the evaluation of the combiner we can filter 
the secret keys by first computing the verification transcript of the combiner on input $(x, y, \pi)$, $\mathcal T$
which then, per our previous argument holds all queryies made to relevant vertices. We then iterate over all candidate secret keys 
$\SK$, any query made by oracle $\O_0$ is trivially simulated, every query made by $\O_{VRF}$ in $\mathcal T$ can honestly be simulated
and every query not in the transcript we iterate over all possible values. If the output over all iterations for a secret key is $y$
(that means that the also in the honest evaluation of the combiner using oracle queries the secret key evaluates to this output), we keep it.
If for some choice of iterated guess the output is not $y$ then we know that the secret key is not the real secret key and we drop it.
We will show that after $2s(\lambda)$ we end up with a candidate set where on a new random $x^*$ the probability is 
significant to sample a secret key which evaluates $x$ to the same output as it is under the honest secret key.
An adversary can leverage this to break pseudorandomness.

\paragraph{Proof of \Cref{lemma:warmup}.}


TODO

Moving on we extend this result to exclude general robust $(n; \lfloor n/2\rfloor)$ black-box combiners.
The proof works very similarly

TODO

\begin{lemma}[Robuts $(n; \lfloor n/2\rfloor)$ Black-Box VRF Combiner]\label{lemma:nhalf}
	There exists no robust $(n; \lfloor n/2\rfloor)$ black-box combiner in the sense of \Cref{def:robust_blackbox_combiner} for the VRF primitive 
	defined in \Cref{def:vrf} for all $n \in \N^+$.
\end{lemma}

The proof of this lemma is very similar to the proof of \Cref{lemma:warmup}, we will therefore move it to the appendix. 
Instead we work on strengthening the result to the following, more general statement.

\begin{theorem}[Robust $(n;n-1)$ Black-Box VRF Combiner]\label{thm:impossibility_wvrf_comb}
	There exists no robust $(n;n-1)$ black-box $\prim{VRF}$-combiner.
\end{theorem}

Extending the proof of \Cref{lemma:warmup} and \Cref{lemma:nhalf} to this stronger notion requires one main insights:
First, repeating the argument before we end up with $n$ oracle queries, of which we can only make $1$ trivial (setting it to $0$ or similar).
The issues is that across oracle calls the individually irrelevant queries may not be jointly irrelevant. That is,
changing the outputs of two individually irrelevant oracle queries across two different candidate schemes may result in 
a different output, making simulation non-trivial. 

\paragraph{Proof of \Cref{thm:impossibility_wvrf_comb}.}
TODO

Finally we show that no combiners exist, even if we lower our requirements for the combiner.
That is we only require statistical correctness, statistical unique provability and weak pseudorandomness.
Even if at least $n-1$ candidates implement secure VRFs, satisfying perfect unique provability,
perfect correcntess and strong pseudorandomness. We state the following theorem, the strongest 
of version of the main result.

\begin{theorem}[Strong Robust $(n;n-1)$ Black-Box VRF Combiner]\label{thm:main}
	There exists no robust $(n;n-1)$ black-box combiner for the VF primitive 
	satisfying statistical correctness, statistical unique provability and weak pseudorandomness.
	Even if at least $n-1$ candidate oracles are guaranteed to implement secure VRFs.
\end{theorem}

\color{red}
have only VRF here, put the other stuff into corollary
\color{black}

The proof extends the proof of \Cref{thm:impossibility_wvrf_comb}. The weak pseudorandomness we get for free in the previous proof,
the constructed adversary can just as well play the weak prf security experiment and with with the same advantage.
To account for the statistical notions of perfect correctness and unique provability we generalize the concept of irrelevance 
to a probabilistic notion called \textit{overwhelming irrelevance}. We say that a node is \textit{overwhelmingly irrelevant}
if the probability, over sampled secret keys $\sk \sample \alg{Gen}{}$ and uniformly random inputs $x \sample \bin^\lambda$
the probability of the vertex being irrelevant is overwhelming. We claim that for overwhelmingly irrelevant vertices the probability of the 
verification algorithm requerying is overwhelming. The proof continues with a similar argument, accounting for the fact that $(1 - \negl)^{\poly}$ is still overwhelming probability.
We provide the entire proof in the appendix. 

We briefly want to comment on the combiner notions on other relaxed version of the VRF primitive.
That is verifiable unpredictable functions (VUF) as defined in \Cref{def:vuf}. We note that for unpredictability 
the same combiner, as given in the proof of \Cref{thm:main} works. Therefore as a corollary to \Cref{thm:main}
we can formulate a similar result on VUFs instead of wVRFs. Since VRF candidates are automatically VUF candidates,
this in particular excludes the possibility of a standard VUF combiner.

\begin{corollary}[Strong Robust $(n;n-1)$ Black-Box VUF Combiner]\label{cor:vuf_comb}
	
\end{corollary}

\section{Weak Robust Black-Box VRF-Combiner}



\begin{comment}
\paragraph{Intuition.}
We need to show that for every $(n;n-1)$ black-box $\prim{VRF}$-combiner $\BBComb{VRF}$
there exists a tuple $\mathbf O = (\O_1, \ldots, \O_n)$ of VF oracles of which at least $n-1$
implement a secure VRF, the combiner $\BBComb{VRF}[\mathbf O]$ doesn't implement a secure VRF.
By equivalence of robust black-box combiner and robust black-box DAG combiners in \Cref{thm:generality_dag_comb}
we instead show the statement for robust black-box DAG combiners. 
For that we fix an arbitrary robust black-box DAG combiner $\BBComb{VRF}$ with the DAGs 
$G_{gen}, G_{eval}, G_{vfy}$ with sources $s_{gen}, s_{eval}, s_{vfy}$, sinks
$t_{gen}, t_{eval}, t_{vfy}$, colorings $c_{gen}, c_{eval}, c_{vfy}$ and ppt. algorithms $\{f_{v}^{gen}\}_{v \in V(G_{gen})}$,
$\{f_{v}^{eval}\}_{v \in V(G_{eval})}$ and $\{f_{v}^{vfy}\}_{v \in V(G_{vfy})}$ for the algorithms $\alg{Gen}{VRF}, \alg{Eval}{VRF}$ and $\alg{Vfy}{VRF}$
respectively. We fix a topological ordering of vertices $v_1, \ldots, v_q$ in $G_{eval}$
and iterate over them. Since we assume verification and evaluation algorithms to be deterministic 
the the randomness from queries to the key-generation algorithm cannot influence the evaluation 
or verification output, therefore those queries can replaced with a hard-coded key.
Additionally since all oracle queries to eval and verification are determiniistic we assume that 
for every vertex the query input for each oracle is unique (since we can just connect an edge 
to the first query, in topolofical ordering, and use this instead of requerying).
We finally make the assumption that $\SK = \bin^{s(\lambda)}$ for $s \in \poly$.

Let $v_1, \ldots, v_n$ be a topological ordering of the vertices in $G_{eval}$. We iterate over the vertices 
in the order $v_1, \ldots, v_n$ and distinguish the following cases: Either for all 
uniquely provable and perfectly correct candidate oracles $\O_1, \ldots, \O_n$ and secret keys $\sk$,
and evaluation points $x$, the output of the 
query made corresponding to $v_i$ does not influence the output at all, that is 
if we change the output of the query associated to $v_1$ the output $y$ of the combiner 
remains the same. We call such vertices \textit{irrelevant}. Irrelevant vertices 
can be simulated by an adversary by using any arbitrary value.\todo{need to add argument why 
replacing the output of irrelevant vertices with arbitrary evaluations preserves unique provability}

Since we know that all prior vertices are irrelevant we know that the function associated 
with the vertex $v_i$, $f_{v_i}^{eval}(\sk, x, o_1, \ldots, o_{d^+(v_i)})$, can be replaced by a function 
$\tilde f_{v_i}^{eval}(\sk, x) = f_{v_i}^{eval}(\sk, x, r_1, \ldots, r_{d^+(v_i)})$ where the values
of $r$ are obtained by arbitrary values in the support of the corresponding oracles.
Assume for all uniquely provable and perfectly correct candidates $\mathbf O$ we have that replacing the true $\sk$ in the execution of the evaluation algorithm with another 
$\widehat\sk, \widetilde\sk \sample \BBComb{VRF}.\alg{Gen}{VRF}(1^\lambda)$ the output of the combiner 
doesnt change, except with negligible probability, then this query is also simulatable as the 
adversary can just query the corresponding oracle on $(\sk', x)$ where $\sk' \sample \BBComb{VRF}[\mathbf O].\alg{Gen}{VRF}(1^\lambda)$
and with overwhelming probability the oracle produces the same output. We call those queries 
\textit{public}. Public queries can also easily be simulated.

If we arrive at a non-irrelevant and non-public vertex $v_j$ then we first claim that the verification algorihm 
necessarily needs to requery the exact same query. If that were not the case then we can construct 
VRF secure candidates (except for the oracle corresponding to this non-irrelevant and non-public vertex),
and only let this oracle differ for that exact evaluation. Since at least $n-1$ candidates are VRF secure,
the combiner also needs to be VRF secure, in particular uniquely provable and perfectly correct.
By perfect correctness the two output however $\hat y \neq \tilde y$ both need to verify
in their respective world. On the other hand in the other world they can not verify 
unless they break unique provability of the combiner. Since the two worlds can only be distinguished 
by this single oracle evaluation we conclude that the verification algorithm must rerun this query.

Let $T \subseteq V(G_{eval})$ be the set of vertices that are not irrelevant and not public.
We now construct the following pseudorandomness adversary $\A = (\A_1, \A_2)$. $\A_1$
plays the experimentation phase of the prf security experiment, it computes $2s(\lambda)$
many input-evaluation pairs $\Brace{(x_i, y_i, \pi_i)}_{i \in 2s(\lambda)}$.
Now it iterates over all evaluation pairs $i \in [2s(\lambda)]$ and iteratively filters the set of 
secret keys $\SK$ using its PSPACE oracle $\PSPACE$. The filtering works as follows: 
For every secret key candidate $\sk \in \SK'$ simulate the computation of the evaluation algorithm,
irrelevant and public vertices are easily simulated. For vertices in $T$ follow the evaluation graph 
and once we would make a query we check whether the query $q$ is in the transcript $\mathcal T_i$,
if the immediate query doesn't occur in $\mathcal T_i$ the secret key $\sk$ can be thrown out.
Otherwise if it occurs in $\mathcal T_i$ then we continue with that oracle response. 
If the final output $y$ differs from the obtained $y_i$ then we can also toss it outside,
otherwise the secret key $\sk$ remains a valid candidate.

Note that due to correctness we never discard the correct secret key $\sk^* \in \SK$,
if in every second filtration we drop at most half the secret keys from the previous candidates
then with high probability, on a uniformly random value $x \sample \bin^*$ we win the pesudorandomness game. 
Otherwise if at least half of the previous candidates are discarded we end up with a polynomially bounded set of 
remaining secret keys $\SK^*$ from which we can guess the correct secret key with significant success probability 
and also win the pseudorandomness game. 

\begin{comment}
For each vertex we distinguish two cases:

\textbf{Case 1}: For all instanciations of the combiner $\BBComb{VRF}$ with $n$ 
uniquely provable and perfectly correct (but not necessarily pseudorandom)
oracles $\mathbf O = (\O_1, \ldots, \O_n)$, secret keys $\sk$ in the support of 
$\BBComb{VRF}[\mathbf o]\alg{Gen}{VRF}$ and $x \sample \bin^\lambda$ the output of the oracle $c(v)$ doesn't 
change the image $y$ produced by the combiner in the sense that for any other output 
$z$ in the support of oracle $c(v)$ would result in the same output. 
In that case an adversary can simulate this query by using any possible response.
The same holds if the output of oracle $c(v)$ only depends on the secret key $\sk$
with negligible probability over randomly sampled $\sk \sample \BBComb{VRF}[\mathbf O].\alg{Gen}{VRF}(1^\lambda)$.
Note that since we iterate over the vertices in a topological ordering we can assume that the input 
of the current iterate can be computed by a $\tilde f_v^{eval}(\sk, x)$ with overwhelming probability.
In that case the adversary can simply guess the secret key and simulate the oracle by running 
the candidate with a negligible decrease in advantage.
If that is the case for all vertices from $v_1$ to $v_q$ then an adversary can simply simulate the 
combiner by sampling a random secret key $\sk$ in the support of $\BBComb{VRF}[\mathbf O].\alg{Gen}{VRF}(1^\lambda)$
and compute a ppt. Turing machine $\tilde f_t^{eval}(\sk, x) = 1$ for the right secret key $\sk$ and $x$ with overwhelming probability, i.e. probability $(1 - \negl)^q$ for $q \in \poly$.
Note that $\tilde f_t^{eval}$ is different to $\BBComb{VRF}[\mathbf O].\alg{Eval}{VRF}(\sk, x)$
not querying any oracles. An adversary can now use its PSPACE oracle $\textsf{PSPACE}$
to filter the space of secret keys $\SK$ on a polynomially sized set of input-image pairs $\{(x_i, y_i, \pi_i)\}_{i \in [2s(\lambda)]}$,
which he can obtain in the pseudorandomness game during the experimentation phase,
on the predicate $\tilde f_t^{eval}(\sk, x_i) = (y_i, \pi_i)$. As we will show, after $2s(\lambda)$ many 
passes of filtering every secret key $\sk$ has overwhelming probability of satisfying 
$\tilde f_t^{eval}(\sk, x^*) = \BBComb{VRF}[\mathbf O].\alg{Eval}{VRF}(\sk, x^*)$ for $x^* \sample \bin^\lambda$.
The adversary $\mathcal D$ can use this to distinguish $(x^*, y^*)$ from $(x^*, r)$ and win the 
pseudorandomness game.

\textbf{Case 2}: On the other hand if \textbf{case 1} doesn't hold for all vertices $v_1, \ldots, v_q$ then 
there exist uniquely provable and perfectly correct VF oracles 
$\mathbf O = (\O_1, \ldots, \O_n)$, a secret key $\sk$ and evaluation point $x$ 
such that there exist two distinct outputs $z_0 \neq z_1$ in the support of the oracle 
$c(v)$ that result in two different images $y_0 \neq y_1$ and with significant probability 
the output of the combiner depends on the choice of the secret key to the evaluation algorithm. 
We now claim that the verification algorithm has to repeat this exact oracle call in order to 
guarantee unique provability and perfect correctness (see \Cref{claim:requery}).
Intuitively if that were not the case then since both the output under $z_0$ and $z_1$
need to verify only if theat single mapping of that single oracle is different 
there is no way to accept only one. This means that within the verification algorithm we have some query to 
$\O_{c(v)}(\tilde f_v^{eval}(\sk, x))$ and thus automatically $\tilde f_v^{eval}(\sk, x)$.
However, note that an adversary might not necessarily be able to tell which query corresponds 
to the requery but that is not important. Specifically let $I$ be the set of 
influential vertices, then it must be that all queries made to oracles 
corresponding to influential vertices are re-queried in the verification algorithm 
of the combiner, since all of the considered queries however sufficiently depend on the 
secret key $\sk$ this means that the secret key $\sk$ can be effectively recovered in that 
scenario. Specifically we record polynomially many transcripts of the verification algorithm,
then for every transcript iterate over the simulated evaluation algorithm, if and only if every oracle query 
being made also occurs in the verification transcript and when used as a simulated answer we obtain 
the correct output and proof we preserve the secret key.
At the end of the filtering on the transcripts $\{\mathcal T_i\}_{i \in 2s(\lambda)}$ we obtain a set 
of secret keys where for any secret key $\sk^*$ and a new random $x^*$ the probability of the simulation agreeing with the true secret key $\sk$
is significant. We can just say that after $2s(\lambda)$ either every second query did reduce the key-space
by half (in that case we end up with a polynomially bounded space with the right key inside, done) or after $2s(\lambda)$ much more keys actually 
resulted in an correct output for many random $x$. Sampling another $x$ u.a.r. it should then be highly likely to also agree there.
 
\todo{todo: why order irrelevant, construct an explicit algorithm for secret key recovery in that case, given only 
the exploration phase of the adversary in the pseudorandomness game} 

Since the evaluation algorithm only makes a polynomial number of queries we can simply remap
secure VRF oracles on the polynomial number of secret keys plus all the secret keys
that share a potential public key with those secret keys for all queries to an 
evaluation oracle, and for every verification oracle we can do the same we can instead obtain 
$\O_1^*, \ldots, \O_n^*$ that are VRF secure and behave identically in this one relevant execution
of the combiner.

\begin{comment}
We claim that in the verification algorithm 
there must exist an oracle call to $c(v)$ on the same input. If that were not the case then 
both the verification of $y_0$ and $y_1$ needs to be accepted in $\BBComb{VRF}[\tilde{\mathbf O}]$
where $\tilde{\mathbf O} = (\O_1, \ldots, \O_{i}[z_0], \ldots, \O_n)$ since by perfect correctness 
both the verification of $y_0$ in $\BBComb{VRF}[\tilde{\mathbf O}]$ and $y_1$ in 
$\BBComb{VRF}[\hat{\mathbf O}]$ must succeed by perfect correctness 
and $\BBComb{VRF}[\hat{\mathbf O}]$ can only be distinguished from $\BBComb{VRF}[\hat{\mathbf O}]$
by querying $\O_i(z_1)$ (as this is the only different evaluation point).

If the response of the oracle query to $c(v)$ is now independent of $x$
except with negligible probability over uniformly randomly sampled $x \sample \bin^\lambda$
then an adversary can also simulate that output since during the experimentation phase as he could simply
run the verification algorithm himself and to obtain the query output \todo{how does he know, it could be unknown which oracle?}
to simulate the query later.

In any case that the oracle $c(v)$ becomes simulatable we simply move on. If we reach the sink node 
$t_{eval}$ then there exists a ppt. Turing machine $f_t^{eval}(\sk, x)$ that makes no oracle queries and 
with overwhelming probability agrees with $\BBComb{VRF}[\mathbf O].\alg{Eval}{VRF}(\sk, x)$.
In the experimentation phase the adversary can compute a polynomial number of distinct tuples $(x_i, y_i, \pi_i)$
such that $f_t^{eval}(\sk, x_i) = (y_i, \pi_i)$ which we can use to filter the secret key space $\SK$
until for any remaining secret key $\sk' \in \SK'$ the probability of a uniformly random secret key agreeing
with the true secret key $\sk$ is non-negligible. The adversary can then simply guess the secret key $\sk$
and win the pseudorandomness game with non-negligible advantage.

Since we assume the robust black-box DAG combiner to be pseudorandom we 
can exclude this case. Therefore there must exist oracles with non-negligible dependence on $\sk$
and influence on the output in the sense that by altering the output to another element in the support 
the output $y$ changes. As already mentioned the same query must be repeated in the verification algorithm.
As a result we can guess the which query corresponds to the re-query and mount the secret key filter attack to 
recover a secret key $\sk$ which matches the overall secret key $\sk$ with non-negligible probability.
Overall the secret key $\sk$ is recovered with probability bounded from below by some $1/q(\lambda)$
for $q(\lambda) \in \poly$ which allows the adversary to break the pseudorandomenss game.

\begin{comment}
To show that there exists no robust $(n;n-1)$ black-box $\prim{VRF}$-combiner $\BBComb{VRF}$
we need to show that every combiner construction $\BBComb{VRF}$ for all oracles $\O_1, \ldots, \O_n$
of which at least $n-1$ implement a secure VRF doesn't implement a secure VRF, i.e. 
is either not perfectly correct according to \Cref{def:correctness}, not uniquely provable 
according to \Cref{def:unique_provability} or not pseudorandom according to \Cref{def:pseudorandomness}.
Since the notions of a robust black-box combiner and a robust black-box DAG combiner we can equivalently
show that there exists no robust black-box DAG combiner. Assume towards a contradiction that there exists 
a robust black-box DAG $\prim{VRF}$-combiner $\BBComb{VRF}$. We now consider the directed acyclic graph 
underlying the computation of the evaluation algorithm $G_{eval} = (V_{eval}, E_{eval})$ and 
the vertices in topological ordering $v_1, \ldots, v_n \in V_{eval}$.
We iterate over those vertices 
and for every vertex $v_i$ we distinguish two cases:

\paragraph{Case 1.} For all VF oracles $\O_1, \ldots, \O_n$ which are uniquely provable 
according to \Cref{def:unique_provability} and perfectly correct according to \Cref{def:correctness}, all secret keys $\sk \in \SK$, all
evaluation points $x \in \bin^\lambda$
and all possible outputs $\tilde\z,\hat\z$ in the range of $c(v_i)$ we have that the $\BBComb{VRF}[\O_1, \ldots, \O_n][v_i \mapsto \tilde\z].\alg{Eval}{VRF}^y(\sk, x)$ 
is equal to $\BBComb{VRF}[\O_1, \ldots, \O_n][v_i \mapsto \hat\z].\alg{Eval}{VRF}^y(\sk, x)$\todo{introduce notation}. In that case for any instanciation of the combiner 
with $n$ uniquely provable oracles the combiner simply ignores the response of the query made at $v_i$ 
and an efficient adversary could simply simulate it with the response $0$. In case that this holds for all 
vertices $v_i \in V_{eval}$ we have that for all uniquely provable VF oracles $\O_1, \ldots, \O_n$
the evaluation of the combiner is equivalent to the evaluation of $f_{t_{eval}}(\sk, x, \mathbf 0)$.
If that is the case then the combiner doesn't rely on its candidate schemes at all and we can use the 
PSPACE oracle $\PSPACE$ to break the combiners pseudorandomness. How that works we explain later.

Alternatively we also consider the case that the output of $f_v(\sk, x, \mathbf 0)$ is indifferent of $x$ except 
with negligible probability, then we can also just guess the output by evaluating with the same $x$ and running the verification 
algorithm which must involve a repeated oracle query call (as we elaborate later as well). 
If it is not indifferent in $x$ with high probability we also distinguish the case that it is indifferent in $\sk$
except with significant probability, in that case we just guess beforehand.

The attack on the combiners pseudorandomness works by iterating over $2s(\lambda)$ many evaluations that 
we can obtain from the oracle in the pseudorandomness game $\{(x_i, (y_i, \pi_i))\}_{i \in [2s(\lambda)]}$ on 
uniformly random inputs $x_i \sample \bin^\lambda$. We then use the PSPACE oracle $\PSPACE$ to iterate over 
all secret keys $\sk \in \SK$ in a random order and check whether $f_{t_{eval}}(\sk, x_i) = (y_i, \_)$ for all $i \in [2s(\lambda)]$.
Note that we always find at least one key $\sk$ since we assume the combiner to be perfectly correct. Let 
$\sk^*$ be the secret key queried to the combiners evaluation algorithm, then we have

\[
\Prob*{f_{t_{eval}}(\sk, x, \mathbf 0) = f_{t_{eval}}(\sk^*, x, \mathbf 0) : x \sample \bin^\lambda} > \frac{1}{2}
\]

This is because if for at least half of the values $(x_i, (y_i, \pi_i))$ the size of the possible 
secret keys is reduced by $\frac{1}{2}$ then the remaining space of secret keys is polynomially bounded
(since we assume that $\SK = \bin^{s(\lambda)}$). If on the other hand at every step on a random $x_i$
at least half the secret keys work and correctly evaluate to the same $y_i$ then the probability of 
randomly sampling a secret key $\sk$ that results in the same output is sufficient to break the game with 
high probability. The adversary $\A = (\A_1, \A_2)$ against pseudorandomness in this case 
is given in \Cref{fig:adversary_pseudorandomness}.

\paragraph{Case 2.} \color{red}salvaged due to modifications above\color{black} Otherwise there exist uniquely provable and perfectly correct VF oracles $\O_1, \ldots, \O_n$,
a secret key $\sk \in \SK$ and an input $x \in \bin^\lambda$ such that the outputs of 
$\BBComb{VRF}[\O_1, \ldots, \O_n][v_i \mapsto \tilde\z].\alg{Eval}{VRF}^y(\sk, x)$ and 
$\BBComb{VRF}[\O_1, \ldots, \O_n][v_i \mapsto \hat\z].\alg{Eval}{VRF}^y(\sk, x)$ are not the same.
Assume that $c(v_i)$ is a key-generation oracle of a candidate, then the combiner must accept both 
\todo{nuance to only include outputs with non-zero probability of being mapped to, but should only be a matter of being more precise} 
by perfect correctness and can therefore not be uniquely provable.
Assume that the oracle call of $c(v_i)$ is made on candidate $\O_j$, then we define $\tilde\O_j, \hat\O_j$ to be the oracle 
$\O_j$ on all evaluations except for the input to $v_i$ which is then queried, which output is set to $\tilde\z$ and $\hat\z$
respectively\todo{include in the DAG combiner that every query is only made once to a single oracle then this also doesnt influence the output anymore for deterministic algorithms}.
Note that neither $\tilde\O_j$ nor $\hat\O_j$ have to be uniquely provable or perfectly correct anymore. 
Given a secure VRF oracle $\O_{\prim{VRF}}$ now construct oracles $\O_1^*, \ldots, \O_{j-1}^*, \O_{j+1}^*, \ldots, \O_n^*$ which agree with
with $\O_{\prim{VRF}}$ on all but a polynomial number of evaluations, agreee with $\O_1, \ldots, \O_{j-1}, \O_{j+1}, \ldots, \O_n$
on all evaluations made in the execution of the evaluation of $\BBComb{VRF}[\O_1, \ldots, \O_{j-1}, \tilde\O_j, \O_{j+1}, \ldots, \O_n].\alg{Eval}{VRF}^y(\sk, x)$
and $\BBComb{VRF}[\O_1, \ldots, \O_{j-1}, \hat\O_j, \O_{j+1}, \ldots, \O_n].\alg{Eval}{VRF}^y(\sk, x)$ and are uniquely provable and perfectly correct.
Since on all but a polynomial number of queries $\O_k^*$ agrees with $\O_{\text{VRF}}$ we can conclude that it must also be pseudorandom 
since the chance of the secret key $\sk$ and challenge $x$ being in this set is negligible.

We give some intuition how this construction works and prove it formally in the formal proof.
Let $\SK^*$ be the set of secret keys on which the evaluation oracle queries to $\O_k$ are queried\todo{kgen nuance},
then we initialize $\O_k^*$ with $\O_k$ on the inputs in $\SK^*$ (which must be polynomially bounded 
by some $p(\lambda) \in \poly$ as the combiner queries all oracles at most a polynomial number of times).
On all secret keys $\SK'$ which don't share a common $\pk$ with a secret key in $\SK^*$ we can simply 
map the result to $\O_{\prim{VRF}}$ since by changing the evaluations we neither break unique provability nor perfect correctness,
we claim that the secret keys in $\SK \setminus (\SK^* \cup \SK')$ can be uniquely mapped as to not violate unique provability and perfect correctness. \todo{change to non-negligible chance of sampling}
The result is that we obtain VRF secure oracles $\O_1^*, \ldots, \O_{j-1}^*, \O_{j+1}^*, \ldots, \O_n^*$
and oracles $\tilde\O_j, \hat\O_j$ that result in different evaluations on $\sk, x$.
Since the combiner is both uniquely provable and perfectly correct in each of the worlds with the oracles 
$(\O_1^*, \ldots, \O_{j-1}^*, \tilde\O_j, \O_{j+1}^*, \ldots, \O_n)$ and
$(\O_1^*, \ldots, \O_{j-1}^*, \hat\O_j, \O_{j+1}^*, \ldots, \O_n)$ the verification algorithm must accept 
only one of the outputs $\hat y, \tilde y$. If the verification algorithm doesnt repeat this exact query
it can not distinguish the two worlds and must either accept both $\tilde y, \hat y$ or none, hence 
we deduce that in the verification algorithm there must be a query to $v_i$. If this query doesnt depend on 
$x$ then the oracle can also be skipped essentially and simulated if the adversary has access to the intermediate results 
in the exploration phase (which he needs due to the requirement for the proof to contain all subinformation for verification). 
\todo{another problem}. If the query doesnt depend on $\sk$ then we can simply simulate it with a randomly sampled $\sk'$.
Hence the query depends on $\sk, x$ which again allows us to recover the secret key $\sk$ by looking at the inputs 
to the verification oracle in the verification combiner and therefore recovering $\{(x_i, (y_i, \pi_i))\}_{i \in 2s(\lambda)}$
and recovering $\sk$.
\end{comment}

\paragraph{Proof (\Cref{thm:impossibility_wvrf_comb})}
Fix a robust $(n;n-1)$ black-box DAG $\prim{VRF}$-combiner $\BBComb{VRF}$ and consider the directed acyclic graph $\G_{eval}$
corresponding to the evaluation algorithm (see \Cref{def:robust_blackbox_dag_combiner}).
Additionally fix a secret key $\sk \in \SK$ and $x \in \bin^\lambda$. We describe the pseudorandomness adversary $\A = (\A_1, \A_2)$ as follows:
$\A_1$ samples $2s(\lambda)$ many evaluation points $x_i \sample \bin^\lambda$ uniformly at random and obtain the evaluation from the PRF oracle $\BBComb{VRF}[\mathbf O].\alg{Eval}{VRF}(\sk, \cdot)$
to construct $\{(x_i, y_i, \pi_i)\}_{i \in [2s(\lambda)]}$. For each tuple $(x_i, y_i, \pi_i)$ the adversary, given the public key $\pk$ computes a transcript of the combiners verification algorithm
$\mathcal T_i$ which contains all oracle queries and responses made as well as the computed functions.
Now additionally he computes \textit{default} values $z_{i,j}^{eval} := \O_i.\alg{Eval}{VRF}(\sk^*_i, x_j)$ for all $i \in [n]$ and $j \in [2s(\lambda)]$ and a $\smash{\sk^* \sample \O_i.\alg{Gen}{VRF}(1^\lambda)}$.
For each $i = 1, \ldots, 2s(\lambda)$ the adversary $\A_1$ computes $\SK_i$ which is defined as: 

\[
	\SK_i := \Brace*{\sk \in \SK_{i-1} : \tilde f(\sk, x_i, \mathcal T_i) = (y_i, \pi_i)} \text { where }\SK_0 := \SK
\]

Where $\tilde f(\sk, x_i)$ runs the evaluation algorithm of the combiner and for every query to an oracle $\O_i$ we simulate the response by
first checking the transcript $\mathcal T_i$ whether an identical query has been made to use, and otherwise use the default value $z_{i,j}^{eval}$.
Since $\tilde f(\sk, x_i, \mathcal T_i)$ doesnt query any oracles, we can compute all of this us§ing the PSPACE complete oracle $\PSPACE$. 
We now argue that following the topological ordering of the oracle queries (vertices in $\G^{eval}$), if a query is not exactly repeated in the 
combiners verification algorithm, and thus available in the transcript, the query itself is \textit{irrelevant} in the sense that if replaced with $z_{i,j}^{eval}$ the evaluation of the combiner doesn't exist.
We formally state this result in the following claim:

\begin{claim}\label{claim:non_irreelvant_queries_repeated}
	For a fixed, hidden secret key $\sk$ and any arbitrary value of $x \in \bin^\lambda$ we have that if an oracle query $\O_{i,j}(\widetilde \sk, \widetilde x)$ is done in the evaluation algorithm 
	$\BBComb{VRF}[\mathbf O].\alg{Eval}{VRF}(\sk, x) = (y,\pi)$ then the query is either exactly repeated in the evaluation of the combiners verification algorithm on the corresponding public key 
	and $y, \pi$: $\BBComb{VRF}[\mathbf O].\alg{Vfy}{VRF}(\pk, x, y, \pi)$ or the value of $y$ stays the same even if the output of $\O_{i,j}(\widetilde \sk, \widetilde x)$ is replaced with any value $z_{i,j}^{eval}$
	that is in the support of $\O_{i,j}$, that is if the query made to $\O_{i,j}$ is \textit{irrelevant}.
\end{claim}

To prove \Cref{claim:non_irrelevant_queries_repeated} we assume towards a contradiction, that the query to $\O_{i,j}(\widetilde \sk, \widetilde x)$ is neither repeated in $\mathcal T_i$,
nor irrelevant to the output $y$. This means that there exist two possible outputs of oracle $\O_{i,j}(\widetilde \sk, \widetilde x)$ that result in different evaluations $y^{(0)} \neq y^{(1)}$. 
This implies that the combiner either is not perfectly correct or uniquely provable and must either accept both $y^{(0)}, y^{(1)}$ or neither. By perfect correctness we have that 
$y^{(0)}$ must correctly verify in the real world and $y^{(1)}$ must correctly verify in a world where the oracle $\O_{i,j}$ output is remapped to the output 
resulting in the evaluation of $y^{(1)}$. Since however the oracle are never queried on the only different remapping of $\O_{i,j}$ both must be accepted in the real world, breaking unique provability. 

Using the result of \Cref{claim:non_irreelvant_queries_repeated} we know that all queries we make that we cannot look up in the verifiers transcript are irrelevant in nature, 
that is we can replace a single one of them with a default value without changing the combiner evaluation.
During evaluation we thus modify the first \textit{unknown} oracle evaluation, i.e. the evaluation not in the transcript.
During evaluation we replace it with a default value, once we reach the next \text{unknown} oracle evaluation.

For all oracles queries of a single oracle in $[n]$ we have that all queries made that are not repeated in 
the verification algorithm are as a group irrelevant for the output,
that means that any of them can be overwitten / reprogrammed. This cannot be done over multiple oracles however

\color{red}
\begin{itemize}
	\item Given a secret key $\sk$ of an oracle $\O_i$ it is still impossible for an efficient algorithm 
		to find a secret key $\sk'$ such that $\sk \sim \sk'$ with non-negligible probability.
		That means that except with non-negligible probability we dont query an oracle on $\sk_0, \sk_1$ 
		such $\sk_0 \sim \sk_1$ in one execution of the evaluation algorithm.
	\item We can modify the output of one candidate to break unique provability, perfect correctness and 
		pseudorandomness. If we assume that the \textit{individually irrelevant} are not \textit{jointly irrelevant}
		then by defaulting the oracle output in topological ordering we necessarily arrive at an oracle query $q$ 
		for which defaulting changes the output. We \textbf{claim} that this output must therefore also be 
		stored in the transcript of the verification, this is because we can simply rewire the previous oracles
		to assume the modified outputs to obtain $\O_1^*, \ldots, \O_n^*$ for which the evaluation of the verification 
		algorithm doesnt change, BUT the evaluation of the combiner now would have a relevant and public query
		not being requeried by the verification.

		Note that for rewiring we need to make sure that the oracles $\O_1^*, \ldots, \O_n^*$
		are still pseudorandom, uniquely provable and correct, i.e. implement secure VRFs.
		For all calls to eval we need to remap all calls to eval under \textit{related} secret keys 
		$\sk^* \sim \sk$ and all corresponding verification queries. For a verification query we need to rewire 
		the corresponding evaluation call, the evaluation calls to all related secret keys $\sk^* \sim \sk$
		and all corresponding verification queries.

		Now assume that we are at an evaluation query to rewire, the only way that the rewiring can impract the 
		output of the combiners evaluation is if there is a \textit{relevant} query to verify with an associated 
		public key on the original oracle output $y$ or if there is a relevant evaluation on a related secret key 
		on that same input. If the latter were the case we could break related key recovery using the combiner.
		If the former is the case then the verification algorithm should come after the evaluation algorithm 
		since otherwise the evaluation combiner algorithm can predict the $y$. If it comes later then
		we can choose this oracle to be the malleable one, modify its output in the eval algorithm, forcing the 
		verification to $0$ at the same time TODO
	\item We claim that if the irrelevant queries are oracle-wide irrelevant, then they are overall also 
		jointly irrelevant. Intuitively this must be the case since, if the oracles were not jointly irrelevant 
		then one oracles irrelevant queries would only be irrelevant for a evaluations of another oracle,
		and relevant for another specific construction (if that were the case and the other construction 
		would be a secure VRF oracle, then relevant the verification combiner needs to repeat those queries). Since neither that 
		oracle, nor the combiner construction, however has knowledge of how that other oracle evaluates 
		this can clearly not be possible. This would indeed only be possible if the combiner could limit those executions 
		to not secure VRF oracles.
\end{itemize}


Now we just need to somehow get rid of the irrelevant oracle calls and replace them with 
a deterministic algorithm to do the key-filtering...
How about a function $f(x,y)$ such that for a fixed $\hat x, \hat y$
we have $f(\hat x, y_0) = f(\hat x, y_1)$ and $f(x_0, \hat y) = f(x_1, \hat y)$
for all $y_0,y_1,x_0,x_1$ but $f(x_0,y_0) \neq f(x_1, y_1)$.

One example is given by $f(x,y) = \max\{x,y\}$.
\color{black}

We now iterate over the oracles in topological ordering.
For every query we precompute a default value which is equivalent to the corresponding oracle query 
on a random input, either $(\sk, x)$ or $(\pk, x, y, \pi)$ where $\sk$ are tested to not be \textit{related}.
Now we iterate over the oracle queries in topological ordering, for each query in the transcript 
we answer with the correct answer, for an answer in the cache we answer with the answer, for every query not in the transcript we default and store the mapping in 
the cache. Assume at step $v_i$ the combiner output changes, as opposed to if we continue computation with the oracles, 
then we know that the query was not in the transcript (otherwise the execution would be the same). 
Since the query is not in the transcript and the query changes the combiner output this must mean that the 
combiner is necessarily not uniquely provable or perfectly correct.

We now aim to show that the combiner, however, must be uniquely provable and perfectly correct,
since we can find VRF secure instanciation of the oracles that align with the evaluation transcript 
and the subsequent honest evaluation. 
Assume towards a contradiction that at least two candidates have no \textit{VRF-secure extension},
that means that they cannot be extended from the existing mappings in a way such that at least 
$n-1$ of them are VRF secure.
That means that there must exist a conflict, we enumerate all possible conflicts at every time step $i$,
assuming the prior modifications left the cnadidates all VRF secure (base case $i = 0$ where $n$ are assumed to be secure)

\begin{itemize}
	\item Assume we encounter an evaluation query with input $(\sk, x)$ and default it to $(y,\pi)$. The only way that this would break 
		unique provability or perfect correctness would be if there exists a related $\sk'$ which outputted $(y' \neq y, \pi)$.
		Assuming that related keys cannot efficiently be found this case only holds with negligible probability. 
		Otherwise there exists a prior verification query which evaluated to $(\pk, x, y, \pi) \mapsto 0$ for $\pk$ associated with $\sk$
		then we simply make this oracle the insecure one. If on the other hand there exists a verification query to $(\pk,)$ \todo{problem}
	\item Assume we encounter a verification query with input $(\pk, x, y, \pi)$ then this can only break a candidates 
		unique provability or perfect correcntess if there exists a verification on $1$ (must be checked)
		on a related public key, we assume that also finding related public keys is hard.
		Or there must exist an evaluation on a related secret key $\sk, x$ that results in $y, \pi$ with $\sk$
		associated to $\pk$. If that also was in the transcript that doesnt work, since we assume all candidates secure then. 
		If the evaluation was defaulted the probability of that happening 
\end{itemize}

\begin{itemize}
	\item A query we defaulted was queried on eval for $(\sk', x)$ to $(y',\pi')$ even though there is another query to $(y, \pi)$ on $(\sk, x)$.
		In that case if $\sk \sim \sk'$ and $y' \neq y$ we have a contradiction to unique provability.
		However note that $\sk \sim \sk'$ requires the combiner to come up with two related secret keys, 
		which we figure is impossible for efficient algorithms.
	\item A query we defaulted was queried on eval for $(\sk', x)$ to $(y',\pi')$ even though there is another query to 
		verify on $(\pk,x,y',\pi') = 0$ or on $(\pk, x,\tilde y \neq y', \tilde\pi) = 1$. We default all verification queries to 
		$0$ and exclude answeres on $y'$ in simulation that we have seen. 
		Either the evaluation is checked and requires evaluation to $y'$ and we answer the verification with $0$,

		Either the verification query evaluated to $1$ by being in the transcript,
		then the output needed to be computed 
\end{itemize}

Note that for all previous queries there is no conflict, either the queries 
are done to evaluation and there we \textbf{assume} that the evaluation algorithm cannot find two related queries,
hance the default answeres just redefine outputs of the equivalence classes of related secret keys. Verification 
queries are answered with $1$ on $(\pk, x, y, \pi)$ if and only if for all replays with a different choice of $y$ of queries 
answered with $y$

If an oracle query cannot be found in the 
transcript or the cache we evaluate it 
We now want to show that if oracle-wide irrelevant queries can be remapped without changing
the combiners evaluation, overall oracles the irrelevant queries are jointly irrelevant as well. 
We proceed with the following argument: By showing that we can remap those outputs in a way such that 
there exists an $n$-tuple of secure VRF oracles that matches this evaluation transcript in topological ordering 
of the irrelevant vertices, then if we encounter a vertex which becomes relevant through that remapping 
this query must necessarily be repeated

---

Assume, for a warm-up that $n = 2$. If we know that at least one candidate is VRF secure then we know that 
during the execution of evaluation of the combiner on $(\sk, x)$ any relevant query made by either oracle 
must be re-queried by the verification algorithm. If that were not the case then simply let this be the malleable oracle,
then produce two outputs, which by perfect correctness both must verify, breaking unique provability.

Since this has to hold for both oracles, just initialize one of the two oracles with the malleable one, this 
execution becomes simulatable entirely. The other execution is given in parts by the transcript of the verification algorithm 
where the entries that are not contained in the transcript can simply be simulated with random evaluations.
\todo{what happens if the query input to one that was previously in the transcript changes as a result 
of the wrong but output-preserving guess of an irrelevant oracle. I think it will be the case that 
for any new transcript produced by guessing the output wont change, and a query done that would change 
the output must then be repeated in verification algorithm. That means we can just guess our way through.}

--- 

\paragraph{Warm up.} Lets first consider the case $n = 2$. 
We consider two cases: Case $1$ is that for all uniquely provable oracles $\O_1, \O_2$,
all secret keys $\sk \sample \alg{Gen}{VRF}$ and all $x \sample \bin^\lambda$ we all relevant oracle queries 
must are repeated in the transcript of the verification algorithm. 
Assume this were not the case, that is, a relevant oracle query for say $\O_1$
would not be re-queried in the verification, then we could obtain two verifiying distinct images $y_0 \neq y_1$
by using correctness but breaking unique provability. Since, in particular, we can make the oracles uniquely 
provable, \todo{unsure}, we solve this case. 

\begin{comment}
Initialize the combiner with an oracle $\mathbf O$ and proceed in topological ordering.
Once we reach a vertex where the oracle query doesnt occur in the transcript we reach vertices 
which are not relevant to the combiner execution. That is for the fixed execution we can replace their output 
with any value adhering to the oracle interface without changing the combiner output. Default this value 
to some $z_{i,j}^{eval}$ and move on. If for every default the probability of changing the output 
is, independently of previous modifications, negligible we win the pseudorandomness game. If there exists a 
vertex $v_i$ for which defaulting to a uniformly random value results in a different output $y_0 \neq y_1$
then the combiner is, again not pseudorandom. It is therefore sufficient to show that for some defaulting 
scheme the oracle candidates remain VRF secure, except with negligible probability.

For that we consider the $i$'th default action and distinguish based on the type of oracle query 
being defaulted. 

\begin{itemize}
	\item If the query is an evaluation query and we didn't already query it on the same oracle 
		and it cannot be found in the verification transcript, then we sample a new secret key $\sk^*$
		using the key generation algorithm of the oracle and evaluate $(\sk^*, x)$.
	\item If the query is a verification query with public key $\pk$, input $x,y,\pi$ then we default it to $0$.
		This only results in a break of perfect correctness if $x,y,\pi$ correspond to e a query $(\sk, x)$
\end{itemize}

First assume that every vertex $v \in V(\G_{eval})$ is \textit{ignored} in the following sense:
compressible
\begin{definition}[Ignored vertices for $\G_{eval}$]\label{def:ignored}
	
\end{definition}

\begin{claim}
	If every vertex $v \in V(\G_{eval})$ is ignored according to \Cref{def:ignored} then there exists an 
	efficient (ppt) adversary $\A = (\A_1, \A_2)$ which breaks psueudorandomness of the combiner.
\end{claim}

Since we arrive at a contradiction the converse of the initial assumption must hold, that is that not every 
vertex $v \in V(\G_{eval})$ is ignored. Then there exists a secret key $\sk \in \SK$, evaluation point 
$x \in \bin^\lambda$, oracles $\mathbf O = (\O_1, \ldots, \O_n)$ and values $\hat z, \tilde z$ such that 
\[
	\hat y = \BBComb{VRF}[\O_1, \ldots, \O_i[v \mapsto \hat z], \ldots, \O_n].\alg{Eval}{VRF}^y(\sk, x)
	\ne
	\BBComb{VRF}[\O_1, \ldots, \O_i[v \mapsto \tilde z], \ldots, \O_n].\alg{Eval}{VRF}^y(\sk, x) = \tilde y
\]
Where $\O_i$ is the candidate oracle associated with $v$. We now construct an oracle $\mathbf O^* = (\O_1^*, \ldots, \O_i, \ldots, \O^*)$

TODO

Suppose, towards a contradiction that there exists a robust $(n;n-1)$ black-box $\prim{VRF}$-combiner $\BBComb{VRF}$
according to \Cref{def:robust_blackbox_combiner_2}. By \Cref{thm:generality_dag_comb} there also exists a robust black-box 
DAG combiner defined in \Cref{def:robust_blackbox_dag_combiner} defined by 

\[
	\Paren*{
		\G_j, c_j, s_j, t_j, \{f_{j,v}\}_{v \in V(\G_j)}
	}_{j \in [3]}
\]

We now consider $\G_2$ which corresponds to the DAG of the combiners evaluation algorithm.


\begin{figure}[H]
	\centering 
	\begin{pcvstack}[space=1em]
		\procedure[linenumbering]
		{Adversary $\mathcal A^{\textsf{PSPACE}, \mathcal O^{prf}_{\BBComb{VRF}[\mathbf O]}}_1(\pk)$}
		{
			\SK' \gets \SK\\
			\pcfor i \in [2s(\lambda)] \pcdo\\ 
			\t x_i \sample \bin^\lambda\\ 
			\t (y_i, \pi_i) \sample \mathcal O^{prf}_{\BBComb{VRF}[\mathbf O]}(x_i)\\ 
			\t \mathcal T_i \gets \textsf{transcript}(\BBComb{VRF}[\mathbf O].\alg{Vfy}{VRF}(\pk, x_1, y_1, \pi_1))\\ 
			\t \pclinecomment{Precompute responses for irrelevant and public vertices:}\\ 
			\t \widetilde \sk \sample \SK\\ 
			\t f_{sim}: v_i \mapsto \O_{c(v_i)}(\tilde f_{v_i}(\sk, i))\\ 
			\t \pclinecomment{Filter the secret keys:}\\ 
			\t \SK' \gets \{\sk \in \SK : \mathcal S^{\textsf{PSPACE}}(\sk, f_{sim}, x_i, y_i, \pi_i, \mathcal T_i) = 1\}\\ 
			\sk^* \sample \SK'\\ 
			\sigma \gets \sk^*\\ 
			\pcreturn \sigma
		}
	\end{pcvstack}
	\begin{pcvstack}[space=1em]
		\procedure[linenumbering]
		{Adversary $\mathcal A^{\mathcal O^{prf}_{\BBComb{VRF}[\mathbf O]}}_2(\pk, x, y, \sigma)$}
		{
			\sk^* \gets \sigma\\
			y^* \gets \BBComb{VRF}[\mathbf O].\alg{Eval}{VRF}(\sk^*, x)\\ 
			\pcreturn y^* = y
		}
	\end{pcvstack}
	\begin{pcvstack}[space=1em]
		\procedure[linenumbering]
		{Simulator $\mathcal S^{\PSPACE}(\sk, f_{sim}^{eval}, x, y, \pi, \mathcal T)$}
		{
			\pclinecomment{All of this is computed in $\O(1)$ due to $\PSPACE$ complete oracle.}\\ 
			v_1, \ldots, v_{|V(G_{eval})|} \gets \tau(V(G_{eval}))\\ 
			\pcfor i \in [|V(G_{eval})|] \pcdo\\ 
			\t \pcif v_i \in \V_{irrelevant} \cup \V_{public} \pcthen\\ 
			\t \t w_i \gets f_{sim}^{eval}(v_i)\\ 
			\t \pcelse \\ 
			\t \t u_i \gets f_{v_i}^{eval}(\sk, x, \{(w_i, v_i) : v_i \in \mathcal N^+_{G_eval}(v)\})\\ 
			\t \t \pcif u_i \notin \mathcal T \text{ at an oracle query to }\O_{c(v)} \pcthen\\ 
			\t \t \t \pcreturn 0 \pccomment{if the query is not repeated in verify alg.}\\ 
			\t \t \pcelse \\ 
			\t \t \t w_i \gets \mathcal T(u_i) \pccomment{use the response of the oracle in verify alg.}\\ 
			\pcreturn w_t = (y, \pi)
		}
	\end{pcvstack}
	\caption{Combiner for VRFs if all candidate schemes are uniquely provable.}
	\label{fig:prf_combiner}
\end{figure}

Having shown that there exists no robust $(n;n-1)$ black-box combiner for VRF we claim 
that the impossibility is a result of the conflict of unique provability and correctness 
with pseudorandomness. We will now strengthen the result slightly by showing that there exists,
in fact no way to combine $n$ VF of which $n-1$ implement secure VRFs such that 
the combiner implements a VF that is statistically correct in the sense of \Cref{def:correctness},
statistically uniquely provable in the sense of \Cref{def:unique_provability} and 
weakly pseudorandom. This excludes combiner constructions with less requirements on the 
security properties of the combiner. 

\begin{theorem}
	Given $n$ VF oracles, of which at least $n-1$ implement a secure VRF in the sense of \Cref{def:vrf},
	there exists no way to combine those schemes in a black-box fashion such that the resulting 
	construction satisfies statistical correctness according to \Cref{def:correctness}, 
	statistical unique provability according to \Cref{def:unique_provability} and 
	weak pseudorandomness in the sense of \Cref{def:pseudorandomness}.	
\end{theorem}

The proof works almost analogous to the proof of \Cref{thm:impossibility_wvrf_comb}, instead of 
considering irrelevant vertices over all inputs $(\sk, x)$ we consider vertices
which are irrelevant over sampled $\sk$ and $x$ which are irrelevant with 
overwhelming probability. Overwhelmingly irrelevant vertices can indeed be simulated again,
since the probability of the input $\sk, x$ being non-irrelevant is bounded by $1 - (1 - \negl)^{\poly}$
which is again a negligible function. Again non-overwhelmingly irrelevant vertices 
must be requeried by the verification algorithm as otherwise, with significant probability,
on a random $(\sk, x)$ input the original output needs to be accepted except with negligible 
probability and there exists a second acceptable input with significant probability. Since every 
non-overwhelmingly irrelevant vertex is requeried by the verification algorithm we can mount the same attack 
as in \Cref{thm:impossibility_wvrf_comb} to break pseudorandomness. Note that exactly the same attack 
also breaks weak pseudorandomness. 
Since the proof is very similar to the one in \Cref{thm:impossibility_wvrf_comb} we skip it here,
instead we provide a full proof in the appendix. \todo{...}

\begin{comment}

	Countless attempts of a proof which all failed on something non-obvious in informal proof :c

\paragraph{Intuition.}
We need to show that for any robust $(n;n-1)$ black-box $\prim{VRF}$-combiner $\BBComb{VRF}$
according to \Cref{def:robust_blackbox_combiner_2} 
there exist $n$ instances of the VF primitive $\O_1, \ldots, \O_n$ of which 
$n-1$ implement secure VRFs such that $\BBComb{VRF}[\O_1, \ldots, \O_n]$
is not a secure VRF. Assume, towards a contradiction, that there exists such a combiner
and $n$ secure VRF oracles $\O_1, \ldots, \O_n$\todo{This seems like a reasonable assumption but Im not sure}.
Using \Cref{thm:generality_dag_comb} we can express $\BBComb{VRF}$ through three DAGs. Specifically we consider the 
DAG of the evaluation algorithm. We iterate over the nodes
in topological ordering. 
First of all note that no oracle call made to the key-generation algorithm is made 
and influences the output (since otherwise the evaluation algorithm would not be deterministic 
and therefore not uniquely provable or perfectly correct)
If for all instanciation of the candidate oracles through uniquely provable and prefectly 
correct VF,
secret keys $\sk$ in the support of the key-generation algorithm and evaluation points
the output of the current node does not influence the output (that is for every output $z \in \mathcal Y \times \mathcal P$ or $z \in \bin$ of the candidate the output doesnt change)
then we continue to the next node. If that is the case for all nodes then 
we can simply instanciate the combiner with any VRF secure candidates, which are in particularly uniquely provable VFs 
and then simulate the combiner by answering every query with $0$, thus breaking pseudorandomness. 
If that is however not the case for every node, say not for node $v$, then there exists $\sk \in \SK, x \in \bin^\lambda$
and uniquely provable oracles $\O_1, \ldots, \O_n$ as well as two different outputs of node $v$ 
$\hat z, \tilde z$ such that the combiner outputs two different values $(\hat y, \hat \pi), (\tilde y, \tilde \pi)$ 
with $\hat y \neq \tilde y$.

First note that no key-generation evaluation can influence the output as 
otherwise the evaluation algorithm wouldn't be deterministic.
Since only a single mapping of the oracle changes the verification algorithm naturally needs 
to re-query this exact oracle to reject the alternative input. Otherwise if that input is not 
queued again in the verification algorithm the transcript of the verification executions 
must be identical and therefore the output must align. Since the verification algorithm 
must make the same oracle call we can now distinguish on whether 
any information of the secret key is used to compute the query input or not.
If information of the secret key $\sk$ and evaluation point $x$ is used then 
we can recover the secret key with high probability or simulate the oracle, moving on.
If neither $\sk$ nor $x$ sufficiently influences the query input the information to compute the output
is public and an adversary may easily break it.

If we end up iterating through all oracle queries and every one turns out to be predictable 
then we arrive at $f_t(\sk, x, \mathbf r)$ where $\mathbf r$ is a vector which we can 
predict with non-negligible probability. When the adversary now plays the randomness game 
then we simply get $\mathbf r$ during the experimentation phase to filter the secret keys
for $2s(\lambda)$ many rounds if $2^{s(\lambda)}$ is an upper bound for the cardinality of $\SK$ on random inputs.
If in at least every second round we drop half the secret keys we end up with a polynomially bounded
set $\SK'$ from which the adversary may simply guess the correct secret key, again with non-negligible
probability to break pseudorandomness. \todo{must use perfect correctness}

\color{red}

If we require the full candidates to be perfectly correct then we can continue like this for 
verification queries that evaluate in $0$ because then we can simply toggle it to $1$ 
and have two valid executions with again the verification needing to requery. 
If verification is $1$ we cannot simply toggle it and argue that both must be accepted in their 
world because for verify $0$ the combiner may not need to be secure as the candidate may not need to be 
perfectly correct. To check that this output was correctly computed the verification algorithm 
either needs to recompute it for 

For evaluation queries we remap it but also need to then have the verification of the oracle 
accept both evaluations (the previous and the other). If that additional change makes the combiner output 
equivalent again then we cannot apply the argument directly, then however the combiner later (!)
queries the verification oracle that was $0$ before and is toggled to $1$ now \todo{this is highly sketchy 
since then the combiner construction would already test for that specifically hard coded difference somewhere}
otherwise the same argument can be made since the merged oracle $\hat\O_i$ and $\tilde\O_i$
can only be distinguished by evaluating it on that specific query. 

Consider all queries made by the verification oracles in the two worlds $\V_1, \V_2$,
we construct an oracle $\O_i^*$ that interpolates between $\hat\O_1$ and $\tilde\O_1$.
First $\O_i^*$ can simply assume all the mappings in $\V_1 \triangle \V_2$
and $\V_1 \cap \V_2$ if $\hat \O(q) = \tilde \O(q)$. Let $q$ be any query in 
$\V_1 \cap \V_2$ with $\hat\O(q) \neq \tilde\O(q)$ if $q$ is a query to verification 
of oracle $\O_i$ if $q$ is a verification query 

If $c(v) = 1$ then done.

If $c(v) = 2$ then if verification happens using verify on a \textit{related} public key $\pk, x, z, \pi$
for $x,z$ the input output but no re-query then the query in correct verification must be $1$ since we assume the candidates 
to be unique. We simply remap $\pk, x, \tilde z, \pi = 1$ and break unique provability.
If we requery on a related secret key and the same input $\sk, x$ then we just move this related
secret key $\sk'$ to a the map to the other output, breaking unique provability.
If both algorithms are used \todo{unclear}


If $c(v) = 3$ then output is $0$ then the verification oracle can only detect this change 
by means of requery, this means that either secret key is exposed or query input is invariant in 
$x$ or $\sk$ and thus becomes simulatable in every case. If output is $1$ then the combiner can 
detect this either by means of requery, if that is the case then simply toggle it and it will 
accept both $\hat y, \tilde y$. However if the output is $1$ the verification algorithm can also 
check by querying evaluation on a related secret key $\sk',x$ to check whether it is $y$
which if succeeds guarantees that the output is $1$ by pref. correctness. Assuming there was a 
way that for any $\sk, x, y, \pi$ the verification in eval was not $1$, then it would be not perfectly 
correct is there was no check.

Since this change is only detectable from the verifiation algorithm by requerying the same 
query using the fact that the combiner is uniquely provable we know that the combiner makes
that query. The pseudorandomness adversary can get a hold of a polynomial number of direct query inputs 
in the exploration phase by querying its oracle on uniformly random inputs. If the output is constant 
except with negligible probability then it is invariant of $x$ and becomes simulatable with overwhelming probability.
We use the PSPACE oracle to filter the secret keys, doing that we either arrive at an exponentially sized set 
for which each secret key is equivalent form this oracles perspective $\dagger$ 
or a polynomially sized set, at which point we can guess it and win the entire game.
In case of $\dagger$ or invariance in $x$ then we move on to the next, being able to simulate it 
with overwhelming probability.

before constructing the vrf oracle $\O_1^*$ if the query is only checked by 
verification (for eval the pk check and for verify the recomputation over equivalent pk)
then this cannot correspond to $0$ in the vfy execution since 
for eval in eval algorithm we assume correctness so any relevant check relies on $1$
and for verification recomputation means in eval there is evaluation to $0$ which either depends on the 
secret key $\sk$ in some parameter in which case we recover it using the verification algorithm or it doesnt 
at which point the oracle becomes simulatable and we move on. If there only is a verification query with $1$ 
and no additional eval check then we can simply break unique provability by toggling this bit in the new oracle 
to allow both $\hat y, \tilde y$ to be accepted (even if verify output indep.).
We conclude that every query needs to be checked by eval query, for eval queries in eval this 
allows recovery or invariance. 

This only leaves verify in eval to eval in verify. If the query is verify $(\pk, x, y, \pi) = 0$
then we can always toggle it since with eval we cannot detect $0 \to 1$. If $1$ 
then we can not simply toggle it without breaking perfect correctness of a candidate,
however the verification algorithm therefore explicitly needs to check $(\sk, x) ?= y$
for some $\sk \sim \pk$.
then the only possible evaluation query that constrains this is $(\sk, x) ?= y$
\color{black}

Let oracle $\O_i$ denote the oracle which is queried at node $v$ then define $\hat\O_i$ and 
$\tilde\O_i$ to be the oracles that agree with $\O_i$ on every input but reply with $\hat z, \tilde z$
on the input to query $v$ with the verification algorithm being $\O_i$ with the remapping of 
the verification on $(\hat y, \hat \pi), (\tilde y, \tilde \pi)$ to $1$ respectively to preserve 
perfect correctness \color{red}show that this cannot change the combiner evaluation. I think we can still all map to 1 
but not fully sure\color{black}. We will show that we can find oracles $\O_1^*, \ldots, \O_{i-1}^*, \O_{i+1}^*, \ldots, \O_n^*$
that implement secure VRFs but still agree with $\O_1, \ldots, \O_{i-1}, \O_{i+1}, \ldots, \O_n$
on polynomially many points, which is sufficient to have them behave identically within this 
execution of the evaluation algorithm of the combiner. 

We claim and show that the verification algorithm can only distinguish the two worlds by effectively
requeueing the same query made in the evaluation of the eval algorithm. Since that is the case 
we can simply the verification algorithm during the pesudorandomness game exploration phase 
to recover polynomially many inputs to the queried oracle. Now informally if this input 
doesnt correlate with the secret key $\sk$ of the combiner then the oracle query again becomes simulatable 
and we move in in the nodes of the tree, falling back into case 1. If the 
input actually depends on the secret key then we can simply filter the set of secret keys
a polynomial number of times using our pspace oracle (and the fact that all previous inputs become predictable as we iterate through it in topological ordering)
to either reduce the set of possible secret keys to a polynomial set (at which point the problem is solved)
or to an exponential set of secret keys where for every secret key the probability of the output being the same as
for the true secret key is overwhelming, again allowing the continuation of this tactic on further nodes.

, corresponding to candidate oracle calls in 
topological ordering. If for all secret keys $\sk$ and $x$ the oWe need to show that for any robust $(n;n-1)$ black-box $\prim{VRF}$
utput and all possible outputs of the oracle 
call the combiner output remains the same with overwhelming probability then
query is simulatable by our adversary and we replace the oracle query with $1$ and move on.
If we arrive at a query where the chance of differing outputs is non-negligible 
then we show that in the verification algorithm of the combiner 
the same query must be made since otherwise we could rewire this single mapping 
to obtain two different evaluation queries but while keeping the transcript of verification and thus the 
verification verdict the same, which is $1$ by perfect correctness breaking unique provability\footnote{this doesnt work if verify = 1 in eval since remapping makes the candidate not perfectly correct}.

If the verification algorithm however makes this exact query we can obtain the immediate input $(\tilde \sk, \tilde x)$
to the oracle in the expreimentation phase of the pseudorandomness game. If the input $(\tilde \sk, \tilde x)$
sufficiently correlates with the exact secret key then we can recover it using the PSPACE oracle. If the input 
doesn't correlate sufficiently with the query then it becomes simulatable and we move on as the adversary can just 
recompute the result. If we traversed all nodes we have either recovered the secret key or found 
that under all inputs $(\sk, x)$ all queries can be simulated by the adversary we can break pseudorandomness
of the combiner.

we can always find $n$ instances $\O_1, \ldots, \O_n$   
of which at least $n-1$ securely implement the VRF primitive according to \Cref{def:vrf} we can break either 
statistical correctness as defined in \Cref{def:correctness}, unique provability as defined in \Cref{def:unique_provability},
or weak pseudorandomness as defined in \Cref{def:weak_pseudorandomness}. 

Let $\BBComb{VRF}$ be a robust $(n;n-1)$ black-box $\prim{VRF}$-combiner.
Let $\O_2, \ldots, \O_n$ 
be arbitrary but uniquely provable and perfectly correct VF oracles. 
\textbf{Case 1.} If for all uniquely provable and prefectly correct VF $\hat\O_1, \tilde\O_1$
the evaluation algorithms $\BBComb{VRF}[\hat\O_1, \O_2, \ldots, \O_n].\alg{Eval}{VRF}^y = 
\BBComb{VRF}[\tilde\O_1, \O_2, \ldots, \O_n].\alg{Eval}{VRF}^y$ then w.l.o.g. if that 
holds for all $i \in [n]$ then if $\BBComb{VRF}$ is instanciated with $n$ secure VRF oracles 
there exists an efficient prf adversary $\mathcal A$ even though the combiner should be secure 
by assumption. The adversary works by sampling $2s(\lambda)$ many input-output pairs $\{(x_i, y_i, \pi_i)\}_{i \in [2s(\lambda)]}$
and then sampling a uniformly random secret key $\sk$ that satisfies $f_t(\sk, x_i) = (y_i, \pi_i)$
for all $i \in [2s(\lambda)]$. Given a challenge $x^*$ and $y^*$ the adversary then checks whether 
$f_t(\sk, x^*) = (y^*, \pi^*)$ for some $\pi^* \in \P$.

\textbf{Case 2.} If on the contrary there exist uniquely provable and perfectly correct VF oracles 
$\hat\O_1,\tilde\O_1, \O_2, \ldots, \O_n$ such that $\BBComb{VRF}[\hat\O_1, \O_2, \ldots, \O_n].\alg{Eval}{VRF}^y \ne 
\BBComb{VRF}[\tilde\O_1, \O_2, \ldots, \O_n].\alg{Eval}{VRF}^y$ and in particular for some $\sk, x$
we have $\BBComb{VRF}[\hat\O_1, \O_2, \ldots, \O_n].\alg{Eval}{VRF}^y(\sk, x) \ne 
\BBComb{VRF}[\tilde\O_1, \O_2, \ldots, \O_n].\alg{Eval}{VRF}^y(\sk, x)$ then we show that there also exist 
$\hat\O_1, \tilde\O_1, \O_2, \ldots, \O_n$ oracles that are secure VRFs such that agree on all inputs
queried during the execution of the combiner on $(\sk, x)$ input.
We show this by starting from the uniquely provable and perfectly correct VRF oracles and then changing the mapping 
of those oracles on any points not queried during the execution of the combiner on $(\sk, x)$.
That is given the secure VRF oracle $\O_{\prim{VRF}}$ we define $\hat\O_1.\alg{Eval}{VRF}(\sk', x') = \hat\O_1.\alg{Eval}{VRF}(\sk', x')$
if there exists a secret key $\sk^*$ such that $(\sk', \pk)$ and $(\sk^*, \pk)$
are sampleable key-pairs and during the execution of the combiner $\hat\O_1$ is queried on $(\sk^*, x')$. 
We show that we remap the majority of evaluations of the evaluation algorithm and only preserve a polynomial 
number of inputs from $\hat\O_1$ which makes the new $\hat\O_1$ pseudorandom and we also show that the construction is 
uniquely provable and perfectly correct. We now consider the verification algorithm of the combiner 
on oracles $\hat{\mathcal Z}, \tilde{\mathcal Z}$. Since we have perfect correctness during the execution of the verification 
algorithm the combiner cannot query any key-generation algorithms of its candidates that influence the output, 
otherwise we show that the combiner cannot be uniquely provable or perfectly correct. The algorithm can therefore 
only involve evaluation and verification queries. 
Let $\hat\V, \tilde\V, \hat\E, \tilde\E$ denote the set of queries made to oracle $\hat\O_1$ and $\tilde \O_1$
in verification algorithm and set of queries made to oracle $\hat \O_1$ and $\tilde\O_1$ in evaluation algorithm respectively.
We now construct an oracle $\O_1^*$ such that both $(\hat y, \hat\pi)$ and $(\tilde y, \tilde \pi)$
are accepted by the verification algorithm.

Now note that for every query $q$ made in the evaluation algortihm we must make the same query $q$ in the verification 
algorithm if there exists an alternate query response the modifies the output. First of all we note that the 
key-generation algorithm cannot influence the output since the evaluation algorithm is deterministic. 
Then if an eval query $q$ is made in the evaluation algorithm the verification algorithm can only verify the truthful 
computation by recomputing the same query $q$, or making a verification query on the public key, input point, 
evaluation and proof. If only the verification is made we can start with $n$ secure VRF candidates and simply 
define the verification algorithm of this candidate to always output $1$, then the verification algorithm cannot 
check whether the output is correctly computed and therefore must accept both outputs or none, beaking either perfect correctness
or unique provability. If the evaluation algorithm makes a verification query then the verification algorithm 
can only verify this query by requerying the same query or by evaluating to check. If for the output verify = 0 
then we can simply toggle it, making it impossible to detect the change without requerying, if the output is verify = 1
then we simply set verify = 1 and modify the evaluation mapping for the verification algorithm to reject.
Hence the verification algorithm must requery the oracles on all queries that corrlate to the output.
The pseudorandomness adversary can now do large-scale key filtering to compute the secret key and break pseudorandomness.

If on the contrary there is no large enough correlation of the secret key to the oracle queries then the 
pseudorandomness adversary can simply simulate the oracles and therefore we can eliminate this candidate from the list 
of relevant oracles.

Consider the transcript of the evaluation of the verification algorithm for $\BBComb{VRF}[\hat\O_1, \ldots]$ on $(\hat y, \hat\pi)$
and for $\BBComb{VRF}[\tilde\O_1, \ldots]$ on $(\tilde y, \tilde \pi)$. Since the verification algorithm 
needs to be deterministic to not break unique provability and perfect correctness any responses to queries to the key-generation oracles
have to be ignored in the final combiner output. Now consider all queries made with the same input to the same verify and eval oracles
of a candidate that resulted in a different output. If this query was made to a verification oracle then 
one of the verification oracles returned $0$, \color{red}assuming that the same query was not executed during the 
evaluation algorithm \color{black} we can conclude that this output cannot influence the combiner output, 
otherwise we could flip the output of the candidate outputting bit $0$ to bit $1$ and break perfect correctness. 
Hence we only have conflicts with evaluation queries. 

\begin{comment}
We now aim to construct an oracle $\bar\O_1$ such that $\BBComb{VRF}[\bar\O_1, \ldots].\alg{Vfy}{VRF}(\pk,x,\hat y, \hat \pi) = 1$
and $\BBComb{VRF}[\bar\O_1, \ldots].\alg{Vfy}{VRF}(\pk,x,\tilde y, \tilde \pi) = 1$, hence breaking unique provability. We construction $\bar\O_1$
by merging the oracles $\hat\O_1$ and $\tilde\O_1$. For that consider all oracle queries made by $\BBComb{VRF}[\hat\O_1, \ldots].\alg{Vfy}{VRF}$
and $\BBComb{VRF}[\bar\O_1, \ldots].\alg{Vfy}{VRF}$ through the perspective of the DAG combiner. Since all oracles except $\O_1$
are identical we only consider the oracle queries made or a member of $\O_1$. For that let $\hat\V, \tilde\V$
denote the set of queries of the form $(\sk, x)$ made to eval or of the form $(\pk, x, y, \pi)$
made to verify. We will show that queries to the key-generation algorithm can be ignored 
as they either have no influence to the final result or break perfect correctness or unique provability.
Now note that only $\hat\V \triangle \tilde\V$ are queries that can be simulated by $\bar\O_1$ easily since they 
are only queried by one verification algorithm. For all queries in $\hat\V \cap \tilde\V$ that agree on the input of $\hat\O_1$
and $\tilde \O_1$ we also find no conflict since $\bar\O_1$ can simply simulate them with the joint output. 

Only the queries $q \in \hat\V \cap \tilde\V$ with $\hat\O_1(q) \neq \tilde\O_1(q)$ remain.


For $q \in \hat\V \triangle \tilde\V$ the oracle $\O_1^*$ can simply answer hard-code the corresponding 
correct answer. Then for $q \in \hat\V \cap \tilde\V$ we note that either the query $q$ has no influence 
to the output bit $b$ of the combiner or $q \in \hat\E \cap \tilde \E$, this is since if it has influence 
and were not in $\bar\E$ we could simply hardcode $\bar\E$ to force the output such that bit $b = 0$ in $\bar\V$
but since the evaluation of $\bar\E$ remains the same ($q$ not queried there) we reject a valid output. In case $q$ is 
irrelevant we can just peel it off, hence arriving at $q \in \hat\V\cap \tilde\V \cap \hat\E \cap \tilde\E$
and $\hat\O_1(q) \neq \tilde\O_1(q)$. 

Further assume there exists a query with $(\bar\sk, \bar x)$ to evaluation algorithm
$\hat\O_1.\alg{Eval}{VRF}$ that is not either verified with $(\bar\pk, \bar x, \bar y, \bar \pi)$ verification or 
recomputed with $(\bar \sk, \bar x)$ in the verification algorithm, then the output of eval must be indifferent to this query
since otherwise we could remap $\hat\O_1$ at $(\bar\sk, \bar x)$ to generate a different output $y' \neq \bar y$
such that this output verifies, breaking unique provability.

Since this is not the case we have that for each query in the evaluation algorithm $(\bar\sk, \bar x)$
there exists a query of the type $(\bar\sk, \bar x)$ to the evaluation algorithm in verification or 
to of the type $(\bar\pk, \bar x, \bar y, \bar \pi)$ to the verification algorithm or the query doesnt impact the output and can be omitted.
Now if the output of the verification queries dont align then we can switch the one returning $0$ to $1$
to break perfect correctness. Hence we are left with only evaluation queries for some of the VRF evaluation queries to its oracles where the output differs.
Now assume that the change of outputs on $(\bar\sk, \bar x)$ actually results in no change then we can simply choose it arbitrarily. 

We find that there is only a problem if the verification algorithm queries the evaluation oracle $\O_i$
on the same input $(\sk, x)$ for $\hat\O_i, \tilde\O_i$. If remapping this point for $\hat\O_i$
doesnt change the output of eval but the output of verify then we are done and break perfect correctness. 
If remapping this point for 

\newpage

First argue that all key-gen cannot influence bit. Then argue that all checks of verification must output $1$ or 
not be correlated with output bit since otherwise we simply flip it for $\O_i$ and break perfect correctness \footnote{assuming no query to verification during evaluation}.
Since only verify queries that respond to $1$ are correlated with output moving towards $1$ we can simply hardcode verify = 1 for all verification queries of $\O_i$.
Now the only problem are evaluation queries on the same point with different outputs. 

Note that for every evaluation query made in the evaluation algorithm we need to do the corresponding 
query to the evaluation or verification algorithm in the verify algorithm of the combiner. If for all queries 
make the corresponding evaluation query then the evaluation of the evaluation algorithm becomes simulatable.
Otherwise there exists an evaluation query that is not done, in that case if there is a verification query made with output 
$0$ that influences the output then we simply remap this output to $1$, breaking correctness \footnote{assuming that the evaluation algorithm 
doesn't query the verification oracle himself, which we justify.} if it outputs a $1$ then we can simply remap the evaluation oracle to an abritraryily
new point such that the output changes (if it doesnt its not needed during evaluation and we can drop it, recursing on the rest)
where verification outputs $1$ there.

Let $\BBComb{VRF}$ be a robust $(n;n-1)$ black-box
$\prim{VRF}$-combiner, if for any uniquely provable candidates $\O_0, \ldots, \O_n$
we have that $\BBComb{VRF}[\O_0, \O_2, \ldots, \O_n].\alg{Eval}{VRF}^y$ is the same function 
as $\BBComb{VRF}[\O_1, \O_2, \ldots, \O_n].\alg{Eval}{VRF}^y$ then we can initialize the first candidate scheme 
with a secure instance of a VRF and simulate the exact behavior of the evaluation algorithm
by replacing the responses to the evaluation algorithm by the first oracle with the output $0$.
In the search of the remaining oracles we repeat this approach until we are either left with a single uninitialized
candidate $\O_n$ (\textbf{Case 1}) or encounter an index where there exist at least $i+1$ uniquely provable oracles $\hat\O_i,\tilde \O_i, \O_{i+1}, \ldots, \O_n$
such that for at least some secret key $\sk \in \SK$ and evaluation point $x \in \bin^\lambda$
there evaluation of $\BBComb{VRF}[\hat\O_i, \ldots, \O_n].\alg{Eval}{VRF}^y(\sk, x)$ and  
$\BBComb{VRF}[\tilde \O_i, \ldots, \O_n].\alg{Eval}{VRF}^y(\sk, x)$ differ (\textbf{Case 2}).


\paragraph{Case 1.} We initialize the final candidate scheme with the malleable VRF, that is still 
a uniquely provable and perfectly correct VF, that is the VF that always returns $0$ to any query of its evaluation oracle 
and always returns $1$ to any verification query if $y = 0$. In that case we can simulate the output of the combiner 
on any input $\sk, x$ without queriying any candidate scheme on its evaluation oracle by replacing all 
oracle queries in the computation by the responde of the trivial oracle \todo{which simply output $0$ to eval and matches a valid vrf oracle on the rest}.
We reduced the evaluation of the evaluation algorithm of the combiner on input $(\sk, x)$
to the evaluation of a known ppt. Turing machine $\tilde f(\sk, x)$ without any oracle calls.
Here we rely on the fact that we require a robust black-box combiner to derive its security
from the candidate schemes, which it cannot do here anymore since he doesn't have access to any oracle.
The adversary will guess the key as follows, in the initial stage he will sample a polynomial set of random
inputs $x \in \bin^\lambda$ and query the evaluation oracle to obtain the corresponding images $y \in \Y$, 
then he uses his $\PSPACE$ oracle to compute the number of possible secret keys $\sk \in \SK$
using $\tilde f$. In case every additional input-image pair $(x,y)$ reduces the size of the candidate 
secret keys by $\frac{1}{2}$ then there exists a polynomial $p \in \poly$ such that after $p(\lambda)$
reductions the resulting set of secret keys is polynomially bounded \todo{correct?},
at this point the adversary will sample a random secret key and will win the pseudorandomness 
game with significant advantage. In case the candidate space remains exponential size then for every 
random input $x$ the probability of at least half the secret keys not being consistent with $(x,y)$
is negligible which means that the adversary stops and simply attacks on a random secret key found to 
win with significant probability.

To replace the queries to the key-generation oracles in the combiners evaluation algorithm we 
simply query each candidates oracle once and use the output of the key-generation algorithm for the rest of the simulated 
executions (since the runtime of the combiner is bounded by a polynomial in $\lambda$ we know that any queries 
to the key-generation algorithm must be made with security parameter smaller than that value of $\lambda$ which means that for each oracle
we only need to compute a polynomial number of evaluations of the key-generation algorithm). 
Assuming the output of the combiner depends on the output of a verification query to any candidate \todo{unsure how}.
Now we reduced the evaluation of the combiner to the evaluation of a ppt. Turing machine $f_y(\sk, x, \mathbf 0, \mathbf 0)$
for which there needs to exist an algorithm $\tilde f_y(\pk, x, y, \pi)$ that verifies that the computation ofr $f_y(\sk, x, \mathbf 0, \mathbf 0)$
was executed correctly which has to be independent of the output of any oracle query \todo{...}.

\[
	\forall p \in \poly, \Prob*{
		\begin{aligned}
			x \sample \bin^\lambda\\ 
			(\sk, \pk) \sample \BBComb{VRF}[\O_1, \ldots, \O_n]\alg{Gen}{VRF}(1^\lambda)\\ 
			\BBComb{VRF}[\hat\O_1, \ldots, \O_n].\alg{Eval}{VRF}^y(\sk, x)
			\neq 
			\BBComb{VRF}[\tilde\O_1, \ldots, \O_n].\alg{Eval}{VRF}^y(\sk, x)
		\end{aligned}
	} < \frac{1}{p(\lambda)} 
\]

\paragraph{Case 2.} In the second case construct two oracle candidates $\hat\O, \tilde \O$ 
that implements a secure candidate but for the secret key $\sk$
we output $\hat \O_i$ and $\tilde \O_i$ respectively but the same output on 
any other secret key.
Since we initialize both $\hat\O_i$ and $\tilde\O_i$ with a pseudorandom verifiable functions,
everywhere except at the $\sk$ we show that they are overall verifiable random functions and in particular
also pseudorandom. Using a hybrid argument we show that any adversary can distinguish $\hat\O_i$ and 
$\tilde\O_i$ at most with probability $2\e$ if he is given only query access to the 
key-generation, evaluation and verification algorithm. At the same time for any robust black-box combiner 
the verification algorithm can be used to distinguish $\hat\O_i$ and $\tilde\O_i$. Specifically by the case assumption 
we have

\[
	\exists p \in \poly, \Prob*{
		\begin{aligned}
			x \sample \bin^\lambda\\ 
			(\sk, \pk) \sample \BBComb{VRF}[\O_1, \ldots, \O_n]\alg{Gen}{VRF}(1^\lambda)\\ 
			\BBComb{VRF}[\hat\O_1, \ldots, \O_n].\alg{Eval}{VRF}^y(\sk, x)
			\neq 
			\BBComb{VRF}[\tilde\O_1, \ldots, \O_n].\alg{Eval}{VRF}^y(\sk, x)
		\end{aligned}
	} \ge \frac{1}{p(\lambda)} 
\]

Then assume that an adversary can distinguish the world with each one combiner in. 
If he couldn't distinguish the two worlds he would need to accept both inputs, since he is statistically correct 
he accepts one with overwhelming probability and since he is uniquely provable he accepts at most one. 
This means that in both worlds the verification algorithm manages to distinguish the combiners, and therefore 
in particular the considered oracle candidate.

Since we assume that the combiner is uniquely provable and perfectly correct
we find that the verification algorithm returns $1$ if and only if oracle $\hat \O_i$
is used instead of $\tilde \O_i$. This can however only be distinguished by evaluating the oracle 
at this specific key $\sk$ which implies that the verification algorithm must have knowledge of $\sk$
and can therefore directly break pseudorandomness.

Where the key generation algorithm $\alg{Gen}{VRF}$ is indexed by $1$, the evaluation algorithm $\alg{Eval}{VRF}$ is indexed by $2$, 
and the verification algorithm $\alg{Vfy}{VRF}$ is indexed by $3$. We first show that the final algorithm applied in the evaluation 
algorithm cannot depend on the secret key. We claim the following. 
For the first case we consider that $\forall \hat\O_1, \tilde\O_1, \O_2, \ldots, \O_n$ VRF oracles such that the following holds.

\[
	\forall p \in \poly, \Prob*{
		\begin{aligned}
			x &\sample \bin^\lambda\\ 
			(\sk, \pk) &\sample \BBComb{VRF}[\O_1, \ldots, \O_n]\alg{Gen}{VRF}(1^\lambda)\\ 
			\BBComb{VRF}[\hat\O_1, \ldots, \O_n].&\alg{Eval}{VRF}^y(\sk, x)
			\neq 
			\BBComb{VRF}[\tilde\O_1, \ldots, \O_n].\alg{Eval}{VRF}^y(\sk, x)
		\end{aligned}
	} < \frac{1}{p(\lambda)} 
\]

Let $\O_{\prim{VRF}}$ denote a secure VRF oracle. We instanciate the first oracle $\O_1$ with $\O_{\prim{VRF}}$
and inductively apply this argument to the remaining oracles. In case that we end up substituting 
all oracles with $\O_{\prim{VRF}}$ we are done since they can all be simulated by replacing the oracle 
outputs with $0$, preserving the image with overwhelming probability. If that is the case by the DAG combiner 
construction we end up with a ppt. Turing machine $f_t(\sk, x, \mathbf 0)$ which makes no calls to the oracles 
$\O_\prim{VRF}$ and satisfies the following.

\[
\forall p \in \poly, \Prob*{
	\begin{aligned}
		x &\sample \bin^\lambda\\ 
		(\sk, \pk) &\sample \BBComb{VRF}[\O_{\prim{VRF}}, \ldots, \O_{\prim{VRF}}]\alg{Gen}{VRF}(1^\lambda)\\ 
		\BBComb{VRF}[\O_{\prim{VRF}}, \ldots, \O_{\prim{VRF}}].&\alg{Eval}{VRF}^y(\sk, x)
		\neq 
		f_t(\sk, x, \mathbf 0)
	\end{aligned} 
} < \frac{1}{p(\lambda)}
\]

We construct an adversary $\A$ with access to a PSPACE complete oracle $\PSPACE$ that breaks pseudorandomness
of the combiner. For that the combiner first samples a set of pre-images $\X \subseteq \bin^\lambda$
uniformly at random such that $2^{|\X|} \ge |\SK|$\todo{maybe this bound should be stated differently}.
Then evaluates all points in $\X$ using the evaluation oracle $\O_{\text{eval}}$ in prf security game 
\Cref{fig:prf_wprf_up_experiment} to obtain $\Paren*{(x,\O_{\text{eval}}(x)) : x \in \X}$.
Following this he computes a uniformly random representative $\widetilde\sk$ from the set of candidate secret keys $\widetilde\SK$\todo{describe how this goes}
using the $\PSPACE$ oracle. The adversary can then simply recompute the image $y$ of the challenge $x^*$ in the prf security experiment 
using the guessed secret key $\widetilde \sk$. The adversary is outlined in \Cref{fig:adversary_prf_thm}.

\begin{figure}[H]
	\centering 
	\begin{pchstack}[center, space=1em]
		\procedure[linenumbering]
		{Adversary $\A^{\O_{\text{eval}}, \PSPACE}(1^\lambda, \pk, x^*, y_b)$}
		{
			\X \gets \varnothing\\ 
			\pcfor i \in [p(\lambda)] \pcdo\\ 
			\t x \sample \bin^\lambda\\
			\t y \gets \O_{\text{eval}}(x)\\  
			\t \X \gets \X \cup \{(x, y)\}\\ 
			\sk \sample \SK\\ 
			y^* \sample \BBComb{VRF}[\O_0, \ldots, \O_0].\alg{Eval}{VRF}(\sk, x^*)\\ 
			b' \gets y^* = y_b\\ 
			\pcreturn b'
		}
	\end{pchstack}
	\caption{PRF adversary $\A$.}
	\label{fig:adversary_prf_thm}
\end{figure}

Then computes the set $\SK_{\checkmark} = \Paren*{}$
In the other case we have $\exists \hat\O_1, \tilde\O_1, \O_2, \ldots, \O_n$ VRF oracles such that the following holds

\[
	\exists p \in \poly, \Prob*{
		\begin{aligned}
			x &\sample \bin^\lambda\\ 
			(\widehat\sk, \widehat\pk) &\sample \BBComb{VRF}[\widehat\O_1, \ldots, \O_n]\alg{Gen}{VRF}(1^\lambda)\\
			(\widetilde\sk, \widetilde\pk) &\sample \BBComb{VRF}[\widetilde\O_1, \ldots, \O_n]\alg{Gen}{VRF}(1^\lambda)\\  
			\BBComb{VRF}[\hat\O_1, \ldots, \O_n].&\alg{Eval}{VRF}^y(\widehat\sk, x)
			\neq 
			\BBComb{VRF}[\tilde\O_1, \ldots, \O_n].\alg{Eval}{VRF}^y(\widetilde\sk, x)
		\end{aligned}
	} \ge \frac{1}{p(\lambda)} 
\]

Let $\widehat{\mathcal Z_i}, \widetilde{\mathcal Z_i}$ be the oracles that are constructed from $\hat\O_i, \tilde\O_i$ respectively 
by running their respective key-generation algorithm, the respective evaluation algorithm returning the image 
and an empty proof, and the verification algorithm always accepting all images. 
Now let $x \sample \bin^\lambda, (\widehat\sk, \widehat\pk)\sample \BBComb{VRF}[\widehat\O_1, \ldots, \O_n]\alg{Gen}{VRF}(1^\lambda)$
and $(\widetilde\sk, \widetilde\pk) \sample \BBComb{VRF}[\widetilde\O_1, \ldots, \O_n]\alg{Gen}{VRF}(1^\lambda)$
with $(\widehat y, \widehat\pi) \gets \BBComb{VRF}[\widehat\O_1, \ldots, \O_n].\alg{Eval}{VRF}^y(\widehat\sk, x)$
and $(\widetilde y, \widetilde\pi) \gets \BBComb{VRF}[\widetilde\O_1, \ldots, \O_n].\alg{Eval}{VRF}^y(\widehat\sk, x)$\footnote{Note that $\widehat\sk$ and not $\widetilde\sk$.}.
By statistical correctness $(\widehat\pk, x, \widehat y, \widehat\pi)$ must verify by unique provability if 
the tuple $(\widehat\pk, x, \widehat y, \widehat\pi)$ verifies then $(\widetilde\pk, x, \widetilde y, \widetilde \pi)$
must be rejected with significant probability due to $\widehat y \neq \widetilde y$ (which is the case with significant probability) and unique provability 
of the combiner $\BBComb{VRF}[\widehat{\mathcal Z}]$.

Note that $\widehat{\mathcal Z_i}, \widetilde{\mathcal Z_i}$ are note uniquely provable but 
perfectly correct and pseudorandom VF. Let $\widehat{\mathcal Z} = (\O_1, \ldots, \O_{i-1}, \widehat{\mathcal Z_i}, \O_{i+1}, \ldots, \O_n)$
and $\widetilde{\mathcal Z} = (\O_1, \ldots, \O_{i-1}, \widetilde{\mathcal Z_i}, \O_{i+1}, \ldots, \O_n)$ be two oracles
and for uniformly random evaluation point $x \sample \bin^\lambda$ and key-pair $(\sk, \pk) \sample \alg{Gen}(1^\lambda)$\todo{which Generation algorithm} let 
$(\hat y, \hat \pi) = \BBComb{VRF}[\widehat{\mathcal Z}].\alg{Eval}{VRF}(\sk, x)$ and $(\tilde y, \tilde \pi) = \BBComb{VRF}[\widetilde{\mathcal Z}].\alg{Eval}{VRF}(\sk, x)$.
By statistical correctness we know that 

\begin{align*}
	\BBComb{VRF}[\widehat{\mathcal Z}].\alg{Vfy}{VRF}(\pk, x, \hat y, \hat \pi) = 1 \\ 
	\BBComb{VRF}[\widetilde{\mathcal Z}].\alg{Vfy}{VRF}(\pk, x, \tilde y, \tilde \pi) = 1 
\end{align*}

with high probability. At the same time by unique provability we have that, conditioned on the above verification queries 
accepting the verification queries on the other output must reject the input, that is that 

\begin{align*}
	\BBComb{VRF}[\widehat{\mathcal Z}].\alg{Vfy}{VRF}(\pk, x, \tilde y, \hat \pi) = 0 \\ 
	\BBComb{VRF}[\widetilde{\mathcal Z}].\alg{Vfy}{VRF}(\pk, x, \hat y, \tilde \pi) = 0 
\end{align*}

This means that we can use the verification algorithm to distinguish

\[
	\Brace*{
		\tilde y : 
		\begin{aligned}
			x &\sample \bin^\lambda\\
			(\sk, \pk) &\sample \BBComb{VRF}[\widetilde{\mathcal Z}].\alg{Gen}{VRF}(1^\lambda)\\
			(\tilde y, \tilde \pi) &\sample \BBComb{VRF}[\widetilde{\mathcal Z}].\alg{Eval}{VRF}(\sk, x)\\
		\end{aligned}
	} 
	\text{ and }
	\Brace*{
		\hat y : 
		\begin{aligned}
			x &\sample \bin^\lambda\\
			(\sk, \pk) &\sample \BBComb{VRF}[\widehat{\mathcal Z}].\alg{Gen}{VRF}(1^\lambda)\\
			(\hat y, \hat \pi) &\sample \BBComb{VRF}[\widehat{\mathcal Z}].\alg{Eval}{VRF}(\sk, x)\\
		\end{aligned}
	} 
\]

then we show that 
any efficient adversary cannot distinguish $\widehat{\mathcal Z}, \widetilde{\mathcal Z}$ with non-negligible probability.
We show the following claim. 

\begin{claim}\label{claim:indis}
	The oracles $\widehat{\mathcal Z}$ and $\widetilde{\mathcal Z}$ are computationally 
	indistinguishable.
\end{claim}

\begin{proof}[Proof of \Cref{claim:indis}]
	This is a result of the robustness of the combiner. Since $n-1$ candidates securely implement 
	the VRF primitive the combiner must also securely implement weak pseudorandomness. 
	We perform a sequence of hybrid steps. Formally we define the games $H_0, H_1, H_2$. 
	In game $H_0$ the adversary playes the pseudorandomness game against the combiner 
	$\BBComb{VRF}[\widehat{\mathcal Z}]$ with hidden bit $b = 0$. 
	In game $H_1$ the adversary plays the pseudorandomness game with hidden bit $b = 1$
	and in game $H_2$ the adversary plays the pseudorandomness game against the combiner 
	$\BBComb{VRF}[\widetilde{\mathcal Z}]$ with hidden bit $b = 0$.
\end{proof}

Now intuitively since the oracles are computationally indistinguishable we must have that the 
corresponding evaluation algorithms for a random secret key $\sk$ and input $x$ are also indistinguishable.
Since the verification algorithms are 

\newpage

Since $\widehat{\mathcal Z}$ and $\widetilde{\mathcal Z}$ only differ in their $i$'th oracle we 
can reduce the indistinguishability of $\widehat{\mathcal Z}$ and $\widetilde{\mathcal Z}$ to 
the indistinguishability of $\widehat{\mathcal Z_i}$ and $\widetilde{\mathcal Z_i}$ with a simple
reduction using the fact that the evaluation algorithm secret key inputs only of a query to a candidate 
never depends on keys generated from the secret key of another candidate scheme.

We first show that give two tuples of VF oracles $(\O_1, \ldots, \hat\O_i, \ldots, \O_n)$
and $(\O_1, \ldots, \tilde \O_i, \ldots, \O_n)$ where the candidates $\O_1, \ldots, \O_n$ implement secure VRFs according to \Cref{def:vrf}
we have that if for only a single pair $(\sk^*, x^*)$ the evaluation functions of $\hat\O_i, \tilde \O_i$ disagree then 
and for all other $(\sk, x) \in \SK \times \bin^\lambda$ the evaluation algorithms behave independently, 
we have that for any oracle ppt. Turing machine $\mathcal D$ the probability of distinguishing the two oracles is bounded 
by the following expression.

\[
	\abs*{
		\Prob*{\mathcal D^{\O_1, \ldots, \hat\O_i, \ldots, \O_n}(1^\lambda)} - 
		\Prob*{\mathcal D^{\O_1, \ldots, \tilde\O_i, \ldots, \O_n}(1^\lambda)}
	} \le \frac{p(\lambda)}{2^\lambda|\SK|}
\]

Where $p(\lambda) \in \poly$ is the polynomial bounding the number of queries to its oracles
executed by the distinguisher $\mathcal D$.

\begin{claim}\label{claim:1}
\[
	\forall \sk_1, \sk_2 \in \SK, x \in \bin^\lambda, f_{2,t}(\sk_1, x, \y).y = f_{2,t}(\sk_2, x, \y).y
\]
\end{claim}

\begin{proof}[\Cref{claim:1}]
	Assume towards a contradiction that the claim doesn't hold. Then there exists a two candidate scheme $\O_i^0, \O_i^1$,
	such that for all other fixed choices of key-generation candidates $\O_1, \ldots, \O_{i-1}, \O_{i+1}, \ldots, \O_n$
	the result in the secret key-generation 
	is that there exists two secret keys $\sk_1 \neq \sk_2$ 
	and an evaluation point $x \in \bin^\lambda$ we have $f_{2,t}(\sk_1, x, \y).y \neq f_{2,t}(\sk_2, x, \y).y$.

	We claim that in this case there must exist an algorithm $\tilde f_{2,t}(\pk,x,\y, y)$ that returns the accepting
	bit if and only if $y = f_{2,t}(\sk, x, \y).y$. Intuitively if such an algorithm 
	doesn't exist, then there is no way for the combiner to verify that the output of the final function 
	was actually computed correctly. Specifically this means that the verification algorithm must either be not correct 
	or it must be not uniquely provable.
\end{proof}

\begin{claim}\label{claim:3}
	Queries to any candidates verification or key generation algorithm in the combiners evaluation algorithm
	cannot depend on the secret key of any candidate in a way that affects the output.
\end{claim}

\begin{proof}[\Cref{claim:3}]
	Consider an arbitrary oracle call to either a verification or key generation oracle of a candidate.
	Let $\sim_{out}$ denote the equivalence relation on the outputs of this oracle calls such that two outputs 
	are equivalent under $\sim_{out}$ if and only if the corresponding outputs of the evaluation algorithm are different 
	with non-zero probability. Say towards a contradiction that the input to this oracle call depends on the input secret key,
	specifically let $\sim_{in}$ be an equivalence relation on the inputs of the inputs such that two elements are equal if and only 
	if the current output are equivalent under $\sim_{out}$ and finally let $\sk_1 \sim_{sk} \sk_2$ if and only if the resulting $f_v(\sk_1,x,\y) \sim_{in} f_v(\sk_2,x,\y)$.

	Intuitively if we use the secret key during any oracle call to gen or vfy we can not expect the verification algorithm to 
	correclt verify the output. That is if the final output depends on this output in the sense that there exist two different 
	outputs of the oracle call (for different secret keys) that produce different images $y_0 \neq y_1$, then there exists no 
	way for the verification algorithm to verify the oracle response without knowledge of the secret key which means that 
	either both outputs are accepted, in which case we break unique provability, or both outputs are rejected, which means that 
	we break correctness if we choose one instanciation above.
\end{proof}

Towards the goal of proving \Cref{thm:impossibility_wvrf_comb} we show that if we fix $n-1$ uniquely provable VF oracles
$\O_1, \ldots, \O_{i-1}, \O_{i+1}, \ldots, \O_n$ then regardless of the choice of oracle $\O_i$ the evaluation algorithm
of the combiner remains the same. Formally we state the following lemma.

\begin{lemma}\label{lemma:1}
	If we fix $n$ uniquely provable VF oracles $\O_1, \ldots, \O_n$
	that implement a VRF \Cref{def:vrf} then for all $i \in [n]$ with overwhelming probability the output of the combiners
	evaluation algorithm on a secret key sampled from the key-generation algorithm on a uniformly random point $x \in \bin^\lambda$
	doesnt change if we replace all the repsponses of $\O_i$ during the evaluation process with $0$. Formally 

	\[
		\
	\]
\end{lemma}

\begin{proof}[\Cref{lemma:1}]
	
\end{proof}

Using this result we can construct an adversary $\A$ outlined in \Cref{fig:adversary_prf} against the wprf security experiment
depicted in \Cref{fig:prf_wprf_up_experiment} with advantage $\Adv{VRF}{wprf}{\lambda, \A} \ge 1 - \frac{1}{|\Y|}$. 
$\A$ works be simulating an execution of the evaluation algorithm of the combiner by simulating oracle calls to evaluation
oracles with the response $0$ and recomputing the public values for the verification and key-generation queries,
which don't depend on the secret key by $\Cref{claim:3}$. $\A$ simply recomputes the output of the evaluation 
algorithm on the given input to decide whether it is in the random world or not.
We formalize this intermediate result in the following claim.

\begin{claim}
	There exists 
\end{claim}

Clearly $\A$ is an efficient adversary since we require at most a polynomial number of queries 
to the candidate oracles key generation and verification oracles and don't make any query to the 
wprf oracle in the experiment. Additionally we compute the advantage of the adversary by conditioning on the 
hidden bit $b$ in the security experiment:

\begin{align}
	\Adv{VRF}{wprf}{\lambda, \A} &:= 2 \cdot \abs*{\Prob*{\Exp{VRF}{wprd}(\lambda, \A) = 1} - \frac{1}{2}}\\ 
	&= 2 \cdot \abs*{\Prob*{b = b'} - \frac{1}{2}} \\ 
	&= 2 \cdot \abs*{\frac{1}{2} \cdot \Prob*{b' = 0 : b = 0} - \frac{1}{2} \cdot \Paren*{1 - \Prob*{b' = 1 : b = 1}}} \\ 
	&= \abs*{\Prob*{b' = 0 : b = 0} - \Prob*{b' = 0 : b = 1}} \\ 
\end{align}

By definition we find that $\BBComb{VRF}[\O_1, \ldots, \O_n]$ is indeed not $\errtime$-weakly pseudorandom where
$\epsilon \ge 1 - \frac{1}{|\Y|}$ and $\runtime \in \poly$, and hence by definition $\BBComb{VRF}[\O_1, \ldots, \O_n]$
is not weakly pseudorandom and thus we arrive at a contradiction since we assumed that $\BBComb{VRF}$ is a secure 
robust $(n;n-1)$ black-box combiner.

\begin{figure}[H]
	\centering 
	\begin{pchstack}[center, space=1em]
		\procedure[linenumbering]
		{Adversary $\A^{\O_{eval}}(1^\lambda, \pk, x^*, y_b)$}
		{
			\sk \sample \SK\\ 
			y^* \sample \BBComb{VRF}[\O_0, \ldots, \O_0].\alg{Eval}{VRF}(\sk, x^*)\\ 
			b' \gets y^* = y_b\\ 
			\pcreturn b'
		}
	\end{pchstack}
	\caption{wPRF adversary $\A$.}
	\label{fig:adversary_prf}
\end{figure}

\begin{comment}
We assume, towards a contradiction, that there exists a robust $(n,n-1)$ black-box combiner and therefore
also a robust $(n,n-1)$ black-box DAG combiner. Let $\O_1, \ldots, \O_n$ be the candidate schemes and w.l.o.g. assume that 
$\O_1$ is the malleable scheme. We distinguish two cases. If there exist uniquely provable oracles $\O_2, \ldots, \O_n$ and
two different ways for oracle $\O_1$ to answer its queries such that the computation results in two distinct images we 
break unique provability by initializing $\O_1^0, \O_1^1$ which are constructed to exactly answer the given queries 
each of the given manners while agreeing on all other inputs and accepting every verification query. We claim that the 
given images must be both accepted, namely we show that if the verification algorithm could reject exactly one of the 
images, then we can construct an adversary that can recover the secret key and thus break pseudorandomness. Finally we note that for each 
uniquely provable oracle and a secure VRF oracle we can construct an oracle implementing a secure VRF that agrees on all inputs with 
the uniquely provable oracle on a polynomial number of secret keys.

On the other hand if for all uniquely provable oracles $\O_2, \ldots, \O_n$ and all choices of outputs for oracle $\O_1$
there exists only one image that can be produced then this oracle can essentially be considered dead. In that case we can instanciate 
it with one of the secure $n-1$ oracles, but since it is ignored we can simulate it using the oracle that always answeres each query with $0$. 
We consider the vertices in the DAG $\G_{\BBComb{VRF}[\O_1, \ldots, \O_n]\alg{Eval}{VRF}}$ and remove every node with vertex 
label $\O_1,\O_1.\alg{Eval}{VRF}$ from the graph, replacing their output in the algorithms 
of their successors with $0$. We proceed by induction on the remaining $n-1$ oracles. If we only have two oracles left in every recursion step (including the last one) we entered this 
case, then we can, overall, simulate the entire combiner with $n$ instances of the oracle $\O(0)$. What is remaining is 
a combiner of the form $\widetilde{f_t}(\sk, \x) := f_t(\sk, \x, \mathbf 0)$, we now argue that $f_t$ must actually not depend on the secret key $\sk$.
Again towards a contradiction assume that it does, since the verification algorithm needs to check that the output of $\widetilde {f_t}$ was actually 
computed correctly \todo{at this point it seems like we should also be able to conclude with the fact that no reduction to the secure VRF should be possible.}
we know that there must exist an algorithm $g_t$ in the verification algorithm that must take as input $\pk,x,y,\pi$ and return $1$ if $\widetilde{f_t}(\sk,x) = y$
if the combiner is uniquely provable then there exists at most one $y$ such that $g_t(\pk, x, y, \pi) = 1$ over all $y \in \Y$, an unbounded adversary 
may therefore simply enumerate all values in $\Y \times \P$ to determine which value of $y$ is accepted by $g_t$ and therefore break weak pseudorandomness.
Assume that no such function $g_t$ exist on the contrary, then the verification algorithm doesnt have a mean to verify whether $f_t(\sk, \x, \i)$ is correctly evaluated and therefore 
either always rejects or there exists multiple evaluations that are accepted (otherwise we could have access to an oracle $g_t$).

\begin{proof}[\Cref{thm:impossibility_wvrf_comb}]
	Given a robust $(n;n-1)$ black-box combiner $\BBComb{VRF}$ we initialize it with the following oracles 
	$\O_{\text{vrf}}$ and $\O_{\text{zero}}$ where $\O_{\text{vrf}}$ is a secure VRF oracle and 
	$\O_{\text{zero}}$ is the VF oracle with arbitrary key generation algorithm, evaluation algorithm 
	that responds with $0$ to any query and verification algorithm that accepts any input. We claim that the 
	adversary $\A = (\A_0, \A_1)$ outlined in \Cref{fig:adversary_bbcomb} wins the prf security experiment
	with probability  $1 - \frac{1}{|\Y|}$. $\A$ works by iteratively 

	\todo{todo}
	Hence we may assume that there exists a function $\tilde f_{2,t}$ with the property that
	for all secret keys $\sk \in \SK$ we have the equality $\tilde f_{2,t}(x,\y) = f_{2,t}(\sk,x,\y)$.
	We claim that if we fix $n-1$ uniquely provable VF oracles $\O_1, \ldots, \O_{i-1}, \O_{i+1}, \ldots, \O_n$, the evaluation algorithm of the combiner $\BBComb{VRF}$
	is equivalent independent of the choice of the final oracle $\O_i$. Formally this means the following. 

	\[
		\forall \O_1, \ldots, \O_{i-1}, \O_i, \tilde \O_i, \O_{i+1}, \ldots, \O_n \text{ u.p. VF}, 
			\BBComb{VRF}[\O_1, \ldots, \O_{i-1}, \O_i, \O_{i+1}, \ldots, \O_n].\alg{Eval}{VRF} = 
			\BBComb{VRF}[\O_1, \ldots, \O_{i-1}, \tilde \O_i, \O_{i+1}, \ldots, \O_n].\alg{Eval}{VRF}
	\]

	Assume towards a contradiction that this is not the case. Then this means that there exists two uniquely provable VF oracle $\O_i, \tilde \O_i$
	such that for $n-1$ uniquely provable VF oracles $\O_1, \ldots, \O_n$ there exists a secret key $\sk \in \SK$ and a point $x \in \bin^\lambda$
	such that the respective outputs of the evaluation algorithm differ. Formally this means that.

	\[
		\exists 
	\]

	\todo{} 
	We finally claim that the combiner $\BBComb{VRF}$ is not a secure wVRF and therefore  
	
	Suppose, towards a contradiction, that there exists a robust $(n,n-1)$ black-box $\prim{wVRF}$-combiner $\BBComb{wVRF}$ according to 
	\Cref{def:robust_blackbox_combiner} with candidates indexed by numbers in $[n]$. We first show that neither the evaluation algorithm 
	$\BBComb{VRF}.\alg{Eval}{VRF}$ nor the verification algorithm $\BBComb{VRF}.\alg{Vfy}{VRF}$ may depend on 
	query responses of the key-generation algorithm of the underlying candidate schemes.
	
	We distinguish two cases, by first assuming that 
	the following holds about $\BBComb{wVRF}$.

	\[
		\forall \O_2, \ldots, \O_n \text{ up. VFs}, \sk \in \SK, x \in \bin^\lambda, \O_1^0, \O_1^1: \BBComb{wVRF}[\O_1^0, \O_2, \ldots, \O_n].\alg{Eval}{wVRF}^y(\sk, x) = \BBComb{wVRF}[\O_1^1, \O_2, \ldots, \O_n].\alg{Eval}{wVRF}^y(\sk, x)
	\]

	Then we initialize $\O_1$ with a secure wVRF $\O_{\text{secure}}$ and find that then for all $\sk \in \SK, x \in \bin^\lambda$ and
	any VF candidate schemes $\O_2, \ldots, \O_n$ we have the combinre construction $\BBComb{wVRF}[0, \O_2, \ldots, \O_n]$ is a 
	robust $(n-1, n-2)$ black-box $\prim{wVRF}$-combiner. We proceed by induction until we either fall into the other case or reach robust $(1, 0)$
	black-box $\prim{wVRF}$-combiner. For a robust $(1,0)$ black-box $\prim{wVRF}$-combiner we show that there cannot exist such a construction. 

	\begin{claim}\label{claim:no_n_zero_bb_comb}
		There exists no robust $(n,0)$ black-box $\prim{wVRF}$-combiner.
	\end{claim}

	\begin{proof}[\Cref{claim:no_n_zero_bb_comb}]
		Assume, by contradiction that there exists a secure robust $(n,0)$ black-box $\prim{VRF}$-combiner. Then we instanciate the 
		combiner with the oracles returning $0$ to every query, that is we consider $\BBComb{wVRF}[0,\ldots,0]$. By unique provability of the scheme 
		we have that. An given a public key $\pk$ and any random input $x \in \bin$ with an image $y_b \in \Y$ either the output 
		of the evaluation or a uniformly random output. The adversary uses its PSPACE oracle $\PSPACE$ to iterate all values $y'$ and proof $\pi'$
		and simulate the verification algorithm to compute the set 
		
		\[
			\Brace*{y' \in \Y : \exists \pi' \in \P, \BBComb{wVRF}[0,\ldots,0].\alg{Vfy}{wVRF}(\pk,x,y',\pi')}
		\]

		If this set contains more than one value then we break statistical unique provability with non-negligible probability 
		and if this set is unique then we break weak pseudorandomness with high probability.
	\end{proof}

	If the first case doesnt hold then we must have that there exist oracles $\O_2, \ldots, \O_n$ that implement a VF, 
	a secret key $\sk \in \SK, x \in \bin^\lambda$ and two VF oracles $\O_1^0,\O_1^1$ such that the output of the combiner on $x$ and $\sk$ 
	differs if initialized with the candidates $\O_1^0, \O_2, \ldots, \O_n$ and $\O_1^1, \O_2, \ldots, \O_n$. 
	
	\[
		\exists \O_2, \ldots, \O_n \text{ up. VFs}, \sk \in \SK, x \in \bin^\lambda, \O_1^0, \O_1^1: \BBComb{wVRF}[\O_1^0, \O_2, \ldots, \O_n].\alg{Eval}{wVRF}^y(\sk, x) \neq \BBComb{wVRF}[\O_1^1, \O_2, \ldots, \O_n].\alg{Eval}{wVRF}^y(\sk, x)
	\]
	
	First we claim that the oracles $\O_2, \ldots, \O_n$ can be converted to secure wVRF if give a secure wVRF,
	this works by rewriting all outputs except their output on $\sk,x$ queried during the execution of the combiner
	\todo{unsure if that works}. We now construct the oracles 
	$\widetilde O_1^0, \widetilde O_1^1$ to be the unique oracles that agree with $\O_1^0$ and $\O_1^1$ on all secret keys except for 
	all secret keys queried in one execution of the combiner. We show that there exists no verification algorithm that can distinguish the two
	worlds each we have one of the two combiners. Since the construction is uniquely provable the verification algorithm 
	cannot possible accept both images which implies that he has to reject both images. This implies that for $\BBComb{wVRF}[\O_1, \ldots, \O_n]$
	there $\sk \in \SK, x  \in \bin^\lambda$ there exists no provable image.

	\begin{claim}
		Given two oracles $\O_0, \O_1: \bin^\lambda \to \bin$ such that the number of preimages that evaluate to different bits 
		is bounded by a polynomial in $\lambda$. Then any distinguisher $\mathcal D$ that wins the oracle distinguishability game with 
		significant probability can be used to construct an adversary $\A$ that recovers at least one pre-image.
	\end{claim}
\end{proof}
\end{comment}

We want to mention that instead of achieving weak pseudorandomness we can also allow unpredictability instead.
As a corollary we can exclude the possibility of wVRF combiners in the sense of \Cref{def:weak_pseudorandomness} and
VUF combiners in the sense of \Cref{def:vuf}. 

\begin{corollary}[Impossibility of a Robust $(n;n-1)$ Black-Box $\prim{VRF}$-combiner]\label{cor:impossibility_vrf_comb}
	Given $n$ secure VRF candidate schemes $\O_1, \ldots, \O_n$ there exists no robust black-box combiner that 
	implements a secure VRF scheme if even one candidate scheme is not uniquely provable.
\end{corollary}

\begin{corollary}[Impossibility of a Robust $(n;n-1)$ Black-Box $\prim{VUF}$-combiner]\label{cor:impossibility_vuf_comb}
	Given $n$ secure VUF candidate schemes $\O_1, \ldots, \O_n$ there exists no robust black-box combiner that 
	implements a secure VUF scheme if even one candidate scheme is not uniquely provable.
\end{corollary}

Having ruled out the possibility of combining VRFs, VUFs, wVRFs even for statistical notions of correctness 
and unique provability we need to make an even stronger assumption to allow for combiners to exist.
We will assume that all candidates satisfy unique provability. While this might sound like a very 
strong assumption, which it is, we argue that it is not too hard of a restriction on modern 
schemes that pose potential candidate schemes for a combiner. This is because unique provability 
is an information theoretic notion which is typically mathematically proven for good constructions 
such as ECVRF or Lattice VRF.

However, assuming that all candidates satisfy unique provability combiner constructions become very simple.
We present a construction very similar to the simple xor combiner for PRFs that works as follows.
The key generation algorithm is returns a tuple of $n$ secret keys and a tuple of $n$ public keys, where 
the $i$'th entry is an independently sampled key from candidate $i$, the evaluation algorithm 
evaluates $x$ once for each candidate under the corresponding secret key and XORs the results together to 
obtain the image, additionally the proof is simply the transcript of the execution.
The verification algorithm checks whether all oracle calls align by iterating over the queries made in the transcript 
and verifying each with a call to the corresponding verification oracle with the corresponding public key. 
Afterwards we recompute $y$ by xoring the partial outputs of the candidates to verify that we obtain the same image.
To prevent the combiner from being incorrect if a candidate it incorrect we execute a public checking routine 
before evaluating any algorithm. The public checking routine consists of sampling a polynomial number (in $\lambda$)
of secret keys $\sk$ and evaluation points $x$, computing their evaluation $(y, \pi)$
and then verifying with the corresponding public key $\pk$ for each oracle. If one oracle behaves incorrectly it is ignored 
during the execution of the combiner. Allowing the combiner to achieve statistcal correctness.
Note that this combiner already works if only at least one candidate is secure. We state this result 
as a lemma below and show the construction of the combiner in \Cref{fig:prf_combiner}.

\color{red}
\textbf{DOESNT WORK DUE TO UNIQUE PROVABILITY BREAK
when collision is found but not remarked}
Here couldn't the evaluation simply involve a check to verify. If this check fails then a candidate fails,
this is proven to the verification algorithm by providing an incorrect evaluation/verification in the proof.
If the verifier can detect that fault then he accepts the skipping of the candidate,
if he doesn't detect the fault then he cannot accept the skipping and rejects if that flag is set 
or if the other verification doesnt accept.
\color{black}

\begin{figure}[hptb]
	\centering 
	\begin{pcvstack}[space=1em]
		\begin{pchstack}[center, space=1em]
			\procedure[linenumbering]
			{Algorithm $\BBComb{VRF}[\O_1, \O_2].\alg{Gen}{VRF}(1^\lambda)$}
			{
				(\sk_1, \pk_1) \sample \O_1.\alg{Gen}{VRF}(1^\lambda)\\ 
				(\sk_2, \pk_2) \sample \O_2.\alg{Gen}{VRF}(1^\lambda)\\ 
				\pcreturn ((\sk_1, \sk_2), (\pk_1, \pk_2))
			}
			\procedure[linenumbering]
			{Algorithm $\BBComb{VRF}[\O_1, \O_2].\alg{Eval}{VRF}(\sk, x)$}
			{
				(\sk_1, \sk_2) \gets \sk\\ 
				(y_1, \pi_1) \gets \O_1.\alg{Eval}{VRF}(\sk_1, x)\\ 
				(y_2,\pi_2) \gets \O_2.\alg{Eval}{VRF}(\sk_2, x)\\ 
				y \gets y_1 \oplus y_2\\ 
				\pi \gets (y_1, \pi_1, y_2, \pi_2)\\ 
				\pcreturn (y, \pi)
			}
		\end{pchstack}
		\begin{pchstack}[center, space=1em]
		\procedure[linenumbering]
		{Algorithm $\BBComb{VRF}[\O_1, \O_2].\alg{Vfy}{VRF}(\pk, x, y, \pi)$}
		{
			(\pk_1, \pk_2) \gets \pk\\ 
			(y_1, \pi_1, y_2, \pi_2) \gets \pi\\ 
			\pcreq y = y_1 \oplus y_2\\ 
			\pcreq \O_1.\alg{Vfy}{VRF}(\pk_1, x, y_1, \pi_1) = 1\\ 
			\pcreq \O_2.\alg{Vfy}{VRF}(\pk_2, x, y_2, \pi_2) = 1
		}
		\end{pchstack}
	\end{pcvstack}
	\caption{Combiner for VRFs if all candidate schemes are uniquely provable.}
	\label{fig:prf_combiner}
\end{figure}

\begin{lemma}\label{lemma:prf_combiner}
	There exist a robust $(n;1)$ black-box combiner that combines $n$ uniquely provable VF candidates, of which at least 
	one implements a secure VRF to construct a VF that is statistically correct in the sense of \Cref{def:correctness},
	uniquely provable in the sense of \Cref{def:unique_provability} and pseudorandom in the sense of \Cref{def:pseudorandomness}.
\end{lemma}

\paragraph{Intuition.}
The proof of the security works in two stages. First we show that by doing the public correctness check 
we can detect incorrect candidates with overwhelming probability. Then we show that given any adversary $\mathcal A$ against the pseudorandomness of the combined 
scheme we can construct a reduction $\mathcal R$ against the pseudorandomness of any underlying candidate.
The reduction works by simulating all other candidates faithfully, relying on the oracle in the prf security experiment 
to simulate the attacked candidate, and then use $\A$. Unique provability follows from unique provability 
of all underlying candidates.

\paragraph{Proof (\Cref{lemma:prf_combiner})}
\todo{should be easy}

Having shown that general notion of combiners on VRFs even with relaxed requirements on the security 
properties of the combiner are impossibile and the combiners on VRFs with guaranteed unique provability 
are easy we want to strengthen the result by considering combiners for VRFs (assuming all candidates are 
uniquely provable) with advanced additional properties, namely constrainability as in \Cref{def:constrainableVRF}, key-homomorphism as in \Cref{def:keyhom},
aggregability as in \Cref{def:agg} and finally unbiasability as in \Cref{def:unbias}. 

For the notions of constrainability, key-homomorphism and aggregability this question is quckly answered 
as we we show how the above construction of extended PRF combiner 
can be minimally adapted to combine the notions. We will quickly explain how to extend the construction 
of the extended PRF combiner to capture those notions and then we will move on to the more interesting 
notion of unbiasability, for which the constructions become more complex.

We now consider additional security properties for VRFs, aggregatability that allows the aggregation of many 
secret keys or image-proof pairs to allow for efficient batch verification and unbiasability which prevents 
adversaries from skewing the output distribution of a VRF by maliciously selecting the keys.



\paragraph{Robust Black-Box Combiner for Unbiasable VRF.} We present a robust black-box combiner for unbiasable VRF under the assumption that 
all candidate schemes are uniquely provable. We first note that there exists a compiler to reduce unbiasable VRF to VRF 
under computational Diffie-Hellman (CDH) assumption \cite{}, a combiner that works under the CDH assumption
would be to ignore unbiasability of the candidates and combining normal VRFs and then introducing unbiasability
afterwards \todo{is this even a robust black-box combiner since we cannot reduce unbiasability from the combiner to its candidates.}.
Instead of relying on an additional hardness assumption we prefer to use the unbiasability property of the secure candidates.
We introduce a robust $(2;1)$ black-box combiner for unbiasable VRF based on collision resistant hash functions and 
randomness extractors that is heavily inspired by the construction of unbiasable VRF presented in \cite{}.
The construction works by initially splitting the input $x \in \bin^\lambda$ into two parts $x = x_0 \| x_1$ then hashing $x_1\|0$ and 
$x_1\|1$ using a collision resistant hash function. Finally we evaluate each of the candidate on one point $H(x_1\|0)$ and $H(x_1\|1)$ 
before running a randomness extractor to extract randomness from $x_0$ and adding it to $y_0 \oplus y_1$. We formulate the result formally
as a theorem.

\begin{theorem}[Robust $(2;1)$ Black-Box Combiner for Unbiasabile VRF] 
	There exists a uniform robust $(2;1)$ black-box combiner for unbiasable VRF if all candidate schemes are 
	uniquely provable.
\end{theorem}

\paragraph{Robust Black-Box Combiner for Aggregatable VRF.} We first look at key-homomorphism combiners.
Given a key-homomorphic function given an evaluation point $x \in \bin^\lambda$ we show how to publicly check 
for the homomorphism property of each candidate scheme. This is possible due to an algorithm by Blum \todo{...} \cite{}
which introduces an efficient homomorphism check that runs in time independent of the group size. This means that given 
an evaluation point we can check whether the evaluation function is a group homomorphism with respect to the secret-key argument
with high probability in time independent of the size of the secret-key space. We state this claim in form of a lemma and quickly elaborate on the 
homomorphism check, referring to the original result for the complete proof.

\begin{lemma}[Robust $(2;1)$ Black-Box Combiner for Key-homomorphic VRF]
	
\end{lemma}

Next we consider aggregable VRFs. We show that indeed there exists a combiner 

\begin{comment}
\section{Robust Black-Box DAG-Combiners}

Generally a black-box combiner can be any construction that makes calls to the implementation 
oracles of its candidates and uses the results to compute an output. We introduce a more structured 
version of a black-box combiner which models its computation using a vertex-colored directed acyclic graph,
called a black-box DAG combiner. 

\begin{definition}[Robust $(n;k)$ Black-Box DAG $\mathbf P$-Combiner]
	A $(n;k)$-robust black-box DAG $\mathbf P$-combiner $\G$ is a $|\mathbf P|$-tuple of
	vertex-colored directed acyclic graphs $(\G_1, \ldots, \G_{|\mathbf P|})$
	where each implementation function $f_i$ on input $\x$ is defined by $[\![\G_i]\!](\x)$
	where $[\![\G_i]\!](\x)$ is defined as in \Cref{fig:dag_combiner}. And $([\![\G_i]\!])_{i \in [|\mathbf P|]}$
	is a $(n;k)$-robust black-box combiner for $\P$.
\end{definition}

\begin{figure}
	\centering 
	\begin{pchstack}[center, space=1em]
		\procedure[linenumbering]
		{Algorithm $[\![(\G_i, s_i, t_i, c_i)]\!](\x)$}
		{
			w_s \gets \x\\ 
			\pcfor v \in \texttt{bfs}(\G_i, s) \pcdo\\ 
			\t w_v \gets f_v(\x, (y_w)_{w \in \mathcal N^-(v)})\\ 
			\t y_v \gets \mathcal O_{c(v)}(w_v)\\ 
			\pcreturn w_t
		}
		\procedure[linenumbering]
		{Algorithm $\texttt{bfs}(\G, s)$}
		{
			\text{lin } \gets \{s\}\\ 
			q \gets \{s\}\\ 
			\pcwhile q \neq \varnothing\\ 
			\t q' \gets \varnothing\\ 
			\t \pcfor v \in q\\
			\t \t \text{ lin } \gets \text{ lin } \cup \mathcal N^+(v)\\ 
			\t \t \text q' \gets q' \cup \mathcal N^+(v)\\ 
			\t q \gets q'\\ 
			\pcreturn \text{ lin}
		}
	\end{pchstack}
	\caption{Implementation function from vertex-colored DAG. The algorithm \texttt{bfs} returns the vertices in $\G_i$ 
	in bfs-order starting from the source $s$.}
	\label{fig:dag_combiner}
\end{figure}

As a proof of non-existence of a black-box combiner needs to exclude a large class of 
constructions it is favourable to introduce the additional structure of the DAG combiner to 
reason more systematically about the internal structure of the combiner. 
Importantly we will show that the notion of a robust black-box DAG combiner 
doesn't lose generality with respect to a general robust black-box combiner if we 
assume strict polynomial time combiners. This claim is formalized in the 
following theorem.

\begin{theorem}[Generality of a Robust Black-Box DAG Combiner]
	Let $\BBComb{P}$ be a robust $(n;k)$ black-box $\mathbf P$-combiner that runs in 
	strict polynomial time. Then there exists a robust $(n;k)$ black-box DAG $\mathbf P$-combiner 
	$\mathcal G$ that satisfies the following properties.
\end{theorem}

\begin{proof}
	
\end{proof}

When combining VRFs the most natural idea is to combine $n$ VRFs out of which 
at least $k$ satisfy all the properties of a VRF. We show that there cannot exist 
DAG combiners that achieve all security properties if that is the case, which implies 
that there exists also no general black-box combiner that combines $n$ general VRFs, out 
of which $k$ are secure. More than that we show that even weaker notions of VRFs, weak VRFs and 
VUFs cannot be combined in the generic way. 

\begin{theorem}[Impossibiltiy of a robust $(2;1)$ black-box $\textbf{VUF}$-combiner]
	\label{theorem:impossiblity_vuf}
	There exists no robust $(2;1)$ black-box $\prim{VUF}$-combiner. 
\end{theorem}

\begin{proof}[\Cref{theorem:impossibility_vuf}]
	Assume, towards a contradiction, that there indeed exists a robust $(2;1)$
	black-box DAG $\prim{VRF}$-combiner $\G = (\G_{\textsf{Gen}}, \G_{\textsf{Eval}}, \G_{\textsf{Vfy}})$.
	By definition of a robust $(2;1)$ black-box $\prim{VRF}$-combiner we know that the following holds 

	\[
		\exists \text{ p.p.t.m. } \A \text{ against }\BBComb{VRF}[\O_1, \O_2] \implies 
		\exists \text{ p.p.t.m. } \mathcal R_1^\A \text{ against }\O_1 \land 
		\exists \text{ p.p.t.m. } \mathcal R_2^\A \text{ against }\O_2
	\]

	We distinguish the cases: 

	\textbf{Case 1} if the following condition holds 

	\[
		\forall \O_0 \text{ u.p. oracle VRF }, \sk, x, \z_1, \z_1' \; (w_t = w_t')
	\]

	Then we let $\O_0 = \lambda x.0$ be the oracle that responds to every query with the output $0$ 
	and let $\O_1$ be any secure VRF oracle then the combiner needs to be secure. But it is clearly not PRF 
	secure since we claim that the final function $f_t$ being applied to compute the output cannot depend on the secret key. 

	\begin{claim}
		The functions $f_v$ for all $v \in V$ in $\G_{\textsf{Eval}}$ may not depend on the secret key.
	\end{claim}

	\begin{proof}
		\todo[this]
	\end{proof}

	\todo[and also not clear how to circumvent the problem of vfy and kgen calls in the $f_t(...)$ term]

	Since this holds any adversary $\A$ plaing against the combiner in the PRF security experiment 
	can simply recompute the value $f_t(\mathbf 0)$ to arrive at the same output as $y^*$, hence the adversary 
	wins the security experiment in the case $b = 0$ with probability $1$ and in the case $b = 1$
	with probability $1 - \frac{1}{2^\lambda}$ and $\A$ is even efficient.
	
	\textbf{Case 2} if the case condition of \textbf{Case 1} doesn't hold we have 

	\[
		\exists \O_0 \text{ u.p. oracle VRF }, \sk, x, \z_1, \z_1' \; (w_t \neq w_t')
	\]	

	We now show that if there exists a uniquely provable oracle VRF then, given another secure VRF,
	we can rewrite the outputs for a specific secret key to obtain a secure VRF with a slightly reduced 
	security parameter in PRF proprty that agrees with the outputs of $\O_0$ on every point queried for that 
	specific secret key $\sk$. Specifically we use the following claim:

	\begin{claim}
		Given a uniquely provable VRF oracle $\O_{u.p.}$ and a secure VRF oracle $\O_{vrf}$
		and a secret key $\sk$ we can construct a VRF oracle $\O_{vrf}^*$ that is secure and satisfies 

		\[
			\forall x \in \X (\O_{vrf}^*(\sk,x) = \O_{vrf}(\sk, x))
		\]
	\end{claim}

	\begin{proof}
		\todo[]
	\end{proof}

	Using the result from claim we get a VRF secure oracle $\O_0^*$ and values $\sk, x, \z_1, \z_1'$ 
	such that the evaluation of the eval method differs. We can now construct 

	\todo[problem what if the input to the oracles???, this might always be secure if ew hash it together with the label of the vertex]

	two oracles $\O_1, \O_1'$ that produce the outputs $\z_1, \z_1'$ respectively. Since $\O_0^*$ is indeed a secure VRF 
	we the combiner would indeed need to be VRF secure, but as we have shown there must two imges $w_t, w_t'$ that correspond 
	to the same pre-image $x$. This is a violation of the unique provability property of the combiner.

	\todo[why dont we encounter the problem of discarding the oracles and running a secure instance or at least make calls]

	\todo[extend to VUF]
	\todo[extend to $(n;n-1)$]
\end{proof}

\begin{theorem}[Impossibiltiy of a robust $(n;n-1)$ black-box $\textbf{wVRF}$-combiner]
	There exists no robust $(n;k)$ black-box $\prim{wVRF}$-combiner. 
\end{theorem}

\begin{proof}
	
\end{proof}

Since every $\prim{VRF}$ is automatically also a $\prim{wVRF}$ and a $\prim{VUF}$
we a direct corollary is the impossiblity of a robust $(n;n-1)$ black-box $\textbf{VRF}$-combiner.

\begin{corollary}[Impossibiltiy of a robust $(n;n-1)$ black-box $\textbf{VRF}$-combiner]
	There exists no robust $(n;k)$ black-box $\prim{VRF}$-combiner. 
\end{corollary}

Since unique provability is typically defined to be an information theoretic 
property that doesn't rely on a hardness assumption but is guaranteed by the construction. 
As a result the combination of unique provability is less interesting than the 
provability of the pseudo-randomness (a computational property) since we can guarantee 
that it holds mathematically. To relax the requirements to the combiner we first assume that 
all $n$ candidate schemes are uniquely povable but only $k$ are pseudo-random and find that there 
exists an easy construction of a robust $(n; 1)$ black-box $\prim{VRF}$-combiner that consists of 
the XOR combiner for the PRF primitive where the proof consists of all outputs and proofs and the verification, 
the combiner is detailed in \Cref{fig:up_robust_comb}.
 
It is now of interested to analyze which advanced properties of a VRF can be combined 
if we assume that all candidates satisfy unique provability. For instance what about 
key-homomorphism VRFs or unbiasable VRFs. Let's look at robust black-box combiners for 
unbiasable VRFs first. 

\begin{theorem}
	There exists a robust $(2;1)$ black-box combiner for unbiasable VRF if all candidate schemes
	are guaranteed to be uniquely provable.
\end{theorem}

Another interesting property for VRF is the property of aggregability or key-homomorphism. 
Like unique provability key-homomorphism is a perfect property and also similarly to unique provability 
we can exclude the possibility of a robust black-box combiner. This is the case intuitively since the combiner 
cannot determine which candidate is key-homomorphic and can also not \todo[this seems like another lengthy proof]

\begin{theorem}
	There exists no robust $(2;1)$ black-box combiner for key-homomorphic VRF, even if we assume all 
	candidate schemes to be uniquely provable.
\end{theorem}

\subsection{Universal Verifiable Random Functions}

Universal constructions of a cryptographic primitive satisfy all security properties 
under the sole condition that a secure instance of that primitive exists. A prominent example 
of a universal construction is given by Levins universal one-way function \cite{Impagliazzo1989PseudorandomGF}. 
Following the construction of Levin it is sufficient for a primitive to admit a robust $(2;1)$ black-box combiner 
for there to exist a universal construction for that primitive. However this is however not a necessaily necessary condition 
since in the iterated schemes may be combined in a non-black-box fashion. Theorem therefore doesn't exclude 
the existence of a universal construction. 

\newpage
\newpage

Depending on the underlying cryptographic primitive 
additional simplifications can be made to the DAG-Combiner definition without a loss of generality. 

\begin{definition}[Robust $(n;k)$ Black-Box DAG $\mathbf{VRF}$-Combiner]
	
\end{definition}

\newpage

An $(n;k)$-robust black-box combiner $\BBComb{VRF}$ for a primitive $P$ can be thought of 
computation following a directed acyclic graph $\G = (V, E)$ with two designated vertices $s$ and $t$, the sink and source 
respectively and a vertex coloring $c: V \to [n\times a]$ where $a$ is the number of algorithms 
defined by the primitive $P$. The computation done by the combiner is then defined by following the graph in a 
bfs-fashion and making a call to oracle $c(v)$ to compute $\y_v$ whenever we reach vertex $v \in V$ at the point 
$A_v(\x,(\y_w)_{w \in \mathcal N^-(v)})$ where $A_y$ is a p.p.t.m. algorithm. Formally we introduce the notion of 
a $(n;k)$-robust black-box DAG combiner. 

\begin{definition}[$(n;k)$-Robust Black-Box DAG Combiner]
	
\end{definition}

While it certainly seems easier to reason about the class of DAG-combiners 
than it is to reason about general black-box combiner, we will show that the 
notion of a DAG-combiner is actually strong enough to capture all general black-box combiner
as long as we are guaranteed that all candidate schemes run in strict polynomial time 
and are correct without error. This equivalence is captured in the following theorem: 

\begin{theorem}
	... 
\end{theorem}

As a warm-up we show that there exists a $(n;1)$-robust black-box $\Prim{VRF}$-combiner if all 
$n$ candidate schemes are guaranteed to be uniquely provable. The formal proof of the following lemma 
is deferred to the \Cref{ch:appendix}.

\begin{restatable}{lemma}{restateLemma}
	\label{lemma:3.1}
	There exists a $\Prim{VRF}$-combiner $\BBComb{VRF}[\cdot, \cdot]$ under the 
	guarantee that both candidate schemes are uniquely provable.
\end{restatable}

Using this construction a simple inductive argument shows that there exists a $(n;n-1)$-robust black-box combiner for $\Prim{VRF}$ for $n > 2$ we assume, 
by the induction hypothesis, a construction of a $(n-1,n-2)$-robust black-box $\Prim{VRF}$-combiner $\BBComb{VRF}[\cdots]_{n-1}$ to construct a $(n,n-1)$-robust black-box $\Prim{VRF}$-combiner $\BBComb{VRF}[\cdots]_n$
in the following manner  

\[ 
	\BBComb {VRF}[\mathcal O_{\vrf_1}, ..., \mathcal O_{\vrf_n}]_n = \BBComb {VRF}[\mathcal O_{\vrf_n}, \BBComb {VRF}[\mathcal O_{\vrf_1}, \cdots, \mathcal O_{\vrf_{n-1}}]_{n-1}]
\] 

If $\mathcal O_{\vrf_n}$ is the secure candidate then by the construction of $\BBComb{VRF}[\cdot, \cdot]$
the instance of the combiner is secure. If $\mathcal O_{\vrf_n}$ is not secure then one of the candidates 
$\mathcal O_{\vrf_1}, \ldots, \mathcal O_{\vrf_{n-1}}$ must be secure (since we know that at least on is secure)
in which case by the induction hypothesis $\BBComb{VRF}[\mathcal O_1, \ldots, \mathcal O_{n-1}]_{n-1}$ must be secure
and then again by the \Cref{lemma:3.1} the final construction is secure.

On the other hand we show that there exists no $(2;1)$-robust black-box $\Prim{VRF}$-combiner 
if we drop this guarantee, i.e. if one candidate is not uniquely provable. This result 
holds even if we assume all candidate schemes to be pseudorandom. 
The main result of this thesis is the proof of the following theorem.

\begin{theorem}\label{thm:main}
	There exists no $(n;n-1)$-robust black-box $\Prim{VRF}$-combiner.
\end{theorem}

\section{Impossibility of Robust Black-Box VRF-Combiner}

\section{Relaxed Robust Black-Box VRF-Combiners}

\section{Universal VRFs}

\subsection*{Simplified Setting - Writeup}

Assume, towards a contradiction, that there indeed exists a $(2;1)$-robust black-box combiner 
$\BBComb{\VRF}[\cdot, \cdot]$ for the VRF primitive. We distinguish two cases 

\textbf{Case 1} $\BBComb{\VRF}[\cdot, \cdot]$ satisfies 

\[
	\forall \mathcal O_{\vrf_1} \text{ u.p.}, sk, x, \vec z_0, \neq \vec z_1: 
	f_y(sk, x, \vec z_0, \mathcal O_{\vrf_1}(s(sk), f(\vec z_0))) =
	f_y(sk, x, \vec z_1, \mathcal O_{\vrf_1}(s(sk), f(\vec z_1)))
\]

Then we let $\mathcal O_{\vrf_1} = \mathcal O^0$ be the first candidate scheme that 
responds with $0$ to any query and we let $\mathcal O_{\vrf_2} = \mathcal O^\checkmark$
be the second candidate scheme that securely implements an instance of the VRF primitive. 
The outputs of the secure candidate scheme corresponds to the vectors $\vec z_0, \vec z_1$
to which the output of the eval function of the combiner is by case assumption invariant. 
Therefore we can simply replace the output of the secure candidate scheme with the output of 
a VRF that always responds with $0$ to any query and get 

\todo[any calls to the KGen oracle we can simply imitate from the adversary since 
if we were to get two different images then we would break unique provability.]

% \[
% 	\forall sk,x: \alg{Eval}{\BBComb{\VRF}[\mathcal O^0, \mathcal O^\checkmark]}(sk, x) 
% 	= \alg{Eval}{\BBComb{\VRF}[\mathcal O^0, \mathcal O^0]}(sk, x)
% \]

We want to construct an adversary $\mathcal A = (\mathcal A_0, \mathcal A_1)$ against the PRF 
security experiment of the combiner that wins the experiment with significant advantage. 
To this end we first need to show that the following claim holds about the structure of the combiner 

\begin{claim}
	The output of the final function $f_y(sk, x, \vec z_0, \vec z_1)$ must be invariant of the 
	secret key $sk$. Formally there must exist a known function $\widetilde{f_y}(x, \vec z_0, \vec z_1)$
	that is equivalent to the outputs of $f_y(sk, x, \vec z_0, \vec z_1)$ on any input and for any secret key 
	$sk$.
\end{claim}

\begin{proof}
	The verification algorithm needs to verify whether $y = f_y(sk, x, \vec z_0, \vec z_1)$
	or otherwise more than one image will be provable. Since it needs to verify there must exist a function 

	\[
		\widetilde f_y(pk, x, \vec z_0, \vec z_1, \pi)  
	\]

	that can be used to verify this function note that the more information the proof contains 
	the more feasible this function gets.  We can assume that only oracle outputs are part of the proof 
	since the remaining intermediate values can be computed from the oracle outputs. Now since we have asserted 
	that the output of the all $0$-oracle response execution is a valid output and matches the output 
	we the entire proof must consist of only $0$ and a function value $\vartheta$ computed from the secret key. 
\end{proof}

Since we may now assume that the output of the eval method is given by 

% \[
% 	\alg{Eval}{\BBComb{\VRF}[\mathcal O^0, \mathcal O^\checkmark]}(sk, x) = f_y(x, \vec 0, \vec 0)
% \]

The adversary $\mathcal A = (\mathcal A_0, \mathcal A_1)$ defined as in \Cref{fig:adversary1}
can select a random value $x$ in the first phase and then recompute the value of the challenge
using $x$ and the knowledge about $f_y$ to determine the value of $b'$.

\begin{figure}[H]
	\centering 
	\begin{pchstack}[boxed]
		\procedure[linenumbering]
		{Adversary $\mathcal A_0(\lambda, pk)$}
		{
			x \sample \mathcal X(\lambda) \\ 
			\sigma \gets x\\ 
			\pcreturn (x,\sigma)
		}
		\procedure[linenumbering]
		{Adversary $\mathcal A_1(\lambda, pk, \sigma, y)$}
		{
			y' \gets f_y(\sigma, \vec 0, \vec 0)\\ 
			\pcreturn y = y'
		}
	\end{pchstack}
	\caption{Adversary $\mathcal A = (\mathcal A_0, \mathcal A_1)$ against the PRF security 
	experiment of $\BBComb{\VRF}[\cdot, \cdot]$.}
	\label{fig:adversray1}
\end{figure}

\textbf{Case 2} if \textbf{Case 1} doesn't hold then we have 

\[
	\exists \mathcal O_{v_1} \text{ u.p.}, sk, x, \vec z_0 \neq \vec z_1:  
	f_y(sk, x, \vec z_0, \mathcal O_{\vrf_1}(s(sk), f(\vec z_0))) \neq
	f_y(sk, x, \vec z_1, \mathcal O_{\vrf_1}(s(sk), f(\vec z_1)))
\]

Let $\mathcal O_{v_1}$ be the uniquely provable instance that satisfies the above condition
and let $\mathcal O_{v_0}$ be a VRF defined according to \Cref{fig:construction1} that always outputs $0$ 
to any query to % $\alg{Eval}{\mathcal O_{v_0}}$ and $1$ to every query to the verification oracle. 

% \begin{figure}[H]
% 	\centering 
% 	\begin{pchstack}[boxed]
% 		\procedure[linenumbering]
% 		{Oracle $\alg{Eval}{\mathcal O_0}(sk, x)$}
% 		{
% 			\pcreturn 0
% 		}
% 		\procedure[linenumbering]
% 		{Oracle $\alg{Vfy}{\mathcal O_0}(pk, x, y, \pi)$}
% 		{
% 			\pcreturn 1
% 		}
% 	\end{pchstack}
% 	\caption{Malleable instance of the VRF primitive.}
% 	\label{fig:contsruction1}
% \end{figure}

\begin{claim}
	Both the images $y_0 = f_y(sk, x, \vec z_0, \mathcal O_{\vrf_1}(s(sk), f(\vec z_0)))$ and 
	$y_1 = f_y(sk, x, \vec z_1, \mathcal O_{\vrf_1}(s(sk), f(\vec z_1)))$ are accepted by the 
	combiners verification algorithm. By assumption we have $y_0 \neq y_1$ and hence the 
	combiner is not uniquely provable.
\end{claim}

\todo{maybe make the Ov1 also PRF secure by local rewriting}

\begin{proof}
	Both are accepted by the correctness of the scheme and both are distinct by assumption. 
	This can be made stronger to also work if we consider statistically correct verifiable random functions. 
	That is we require that the correctness property holds only with probability $1 - \negl$, 
	in that case we simply need to let the malleable instance sample a random image at any point to 
	prevent the not-correct points being hard-coded into the combiner.
\end{proof}

\newpage

We now show that the restriction on the verification algortihm captures all valid verification algorithms. 
Intuitively if we have a information theoretically correct instance the outputs of the verification
algorithm are fixed by the outputs of the evaluation algorithm already. For that we define 
$p_{\mathcal O}(\lambda, sk)$ for any secret key $sk$ to be the set of public keys that 
has non-zero probability to be sampled together with $sk$, formally:

% \[
% 	p_{\mathcal O}(\lambda, sk) := \Brace*{pk \in \PK \mid \Prob*{(sk, pk) \sample \alg{KGen}{\mathcal O}(1^\lambda)} > 0}
% \]
% 
% \[
% 	s_{\mathcal O}(\lambda, pk) := \Brace*{sk \in \SK \mid \Prob*{(sk, pk) \sample \alg{KGen}{\mathcal O}(1^\lambda)} > 0}
% \]

\subsection{Extending to more General Constructions}

First note that we can replace the intermediate functions $f_i$ with probabilistic polynomial time 
turing machines since we know that for a fixed $x$ and $sk$ the image $y$ must be the same 
for every execution (since otherwise we break the unique provability requirement assuming correctness). 
So in the case that we break the PRF property the adversary can work the same and in the case that 
we show unique provability the collision is also the same. 

Then we note that there are no restrictions on % $\alg{KGen}{\BBComb{\VRF}[\cdot, \cdot]}$. 

Note that we can simply let Evals arguments depend on both $x$ and $sk$... \todo

Then we note that oracle calls to the functions % $\alg{KGen}{\BBComb{\VRF}[\cdot, \cdot]}$
and % $\alg{Vfy}{\BBComb{\VRF}[\cdot, \cdot]}$ that don't involve the secret key can be easily tolerated 
since the adversary $\mathcal A$ against the PRF experiment can simply do the same that it has done before 
and the collision is still present. 

Then we have the case where we additionally make oracle calls to % $\alg{KGen}{\BBComb{\VRF}[\cdot, \cdot]}$
and % $\alg{Vfy}{\BBComb{\VRF}[\cdot, \cdot]}$ where the call uses arguemnts that depend on the secret key. 
In that case we show that either the output is the same for any secret key (which then allows the adversary to 
again simulate this execution by sampling a random key) or that otherwise the verification cannot succeed \todo 

Now we are left with the only restriction being on the canonical verification of the combiner. Assume that during the verification 
a there exists a query that was made during the execution of the evaluation method that is not verified in that manner. Then if that query has no influence on the output 
then the adversary can simply ignore it during the simulation. If it has an influence there exist outputs $y_0 \neq y_1$ that are both possible 
for some oracle responses. Since verification doesn't check that this oracle query was indeed correctly evaluated 
both $y_0$ and $y_1$ are indistinguishable for any verification algorithm and by correctness both must verify. 
Now say that an intermediate result is not recomputed to check whether the output is consistent with what was queried subsequently,
then we can again differentiate between the case that this input is inrrelevant and then the case must indeed not be covered but we also 
don't rely on the fact that this is checked during the above proof. If however the input has relevance then we can 
again find $y_0 \neq y_1$ that both verify and break unique provability. This means that every important link and every important query must 
be checked during verification. In particular the value of the last function must be checked $\widetilde f_y(pk, x, y, \pi)$ (if that is the identity then this 
check is previously done for the oracle, which cannot be irrelevant since then the adversary could easily determine that its not pseudorandom)
and if it is a function then checking this is inevitable. Hence all relevant features need to be checked. 

Now if the verification algorithm makes additional checks then this cannot fix the attack on the pseudo-randomness 
but only on the unique provability. That means to remedy this attack it would need to distinguish $y_0$ and $y_1$
and determine which one was really produced during the execution. Since the verification algorithm 
is however the same for $\mathcal O_{v_{00}}$ and $\mathcal O_{v_{01}}$ this must happen by 
checking the evaluations. We can now however define the oracles to be identical with exception 
for the specific $sk$ that satisfies the below condition 

\[
	\exists \mathcal O_{v_1} \text{ u.p.}, sk, x, \vec z_0 \neq \vec z_1:  
	y_0 := f_y(sk, x, \vec z_0, \mathcal O_{\vrf_1}(s(sk), f(\vec z_0))) \neq
	f_y(sk, x, \vec z_1, \mathcal O_{\vrf_1}(s(sk), f(\vec z_1))) =: y_1
\]

In that case the difference can only be detected if the oracles are queried on the special secret key 
$sk$ which would mean that the verification algorithm would break the oracle inditinguishability game. 
which means that the verification algorithm must be able to efficiently reconstruct the specific secret key $sk$
in which case we can construct an adversary $\mathcal B = (\mathcal B_0, \mathcal B_1)$ that can reconstruct the secret key 
as well and then trivially break the pseudo-randomness.

\newpage 


We first show \Cref{thm:main} in a simplified setting. 
We initially impose the following restrictions onto the combiner and only exclude combiners that 
respect those assumptions. Also we restrict ourselves initially to $(2;1)$-robust black-box $\Prim{VRF}$-combiners.
The simplifying assumptions are

\begin{itemize}
	\item 
\end{itemize}

In the following we show that there exists no $(2;1)$-robust black-box combiner for the VRF primitive.
Initially we make the following simplifying assumptions on the construction of the combiner.

\begin{itemize}
	\item There exists a directed acyclic graph (DAG) $G = (V,E)$ with two designated vertices $s$ for \textit{source} and $t$ for \textit{target} 
		together with a labling function $l: V \to \bin$ such that the execution of the evaluation algorithm is done by traversing $G$ from $s$ to $t$ in a topological ordering 
		and at every vertex $v \in V \setminus \{s,t\}$ computing the intermediate values 

%		\begin{align*}
%			w_v &= f_v\Paren*{sk, x, \SetBuilder{(u,w_u)}{u \in \mathcal N^+(v)}}\\ 
%			(y_v,\pi_v) &= \alg{Eval}{\vrf_{l(v)}}\Paren*{s_v(sk), w_v}
%		\end{align*}
%
%		We define $y_s = x$ and the $y = w_t = f_t(sk,x, \SetBuilder{(u,w_u)}{u \in \mathcal N^+(t)})$. The proof $\pi$ is the tuple containing all 
%		intermediate values $w_v, y_v, \pi_v$:
%
%		\begin{align*}
%			\pi = \Paren*{\Paren*{w_v}_{v \in V}, \Paren*{y_v}_{v \in V}, \Paren*{\pi_v}_{v \in V}}
%		\end{align*}
%	
%		Where 
%
%		\begin{align*}
%			\forall v \in V: &f_v: \SK_{\Comb\VRF} \times \Paren*{\bin^\lambda}^{1+\deg^+(v)} \to \bin^\lambda\\ 
%			\forall v \in V: &s_v: \SK_{\Comb\VRF} \to \SK_{\vrf_{l(v)}}
%		\end{align*}

		Note that the evaluation algorithm of the combiner doesn't make any calls to the functions % $\alg{Gen}{\vrf}$ and $\alg{Vfy}{\vrf}$
		and is deterministic.
	\item The verification algorithm is constructed in the canonical way by first recomputing the values $w_v$ for all $v \in V \setminus \Paren{s}$ by following the graph in a topological ordering, then verifying
		that $y = w_t$ and finally iterating over the vertices and making a query to the respective candidate to verify if $z_v$ is an image of $w_v$ under $s_v(sk)$. If and only if all those conditions hold the verification algorithm 
		of the combiner accepts.
\end{itemize}

And additionally we assume that 

\begin{itemize}
	\item All candidate schemes are correct 
	\item We assume that there exists a secure instance of the $\VRF$ primitive
\end{itemize}

To show that there exists no $\prim{VRF}$-combiner we show that there exists no 
black-box proof of the combiner, i.e. that there exist oracles $\mathcal O_1, \mathcal O_2$
with the VRF implementation functions such that at least one oracle implements a secure VRF 
and such that there exists an adversary that breaks the security of the combiner $\Comb{VRF}[\mathcal O_1, \mathcal O_2]$
with non-negligible advantage. First we show that there indeed exists an oracle $\mathcal O_{\Prim{VRF}}$
that implements a secure instance of $\prim{VRF}$, the result stated in the following lemma 

\begin{lemma}
	In a relativized model there exists an oracle $\mathcal O_{\prim{VRF}}$
	that implements a secure instance of a $\prim{VRF}$.
\end{lemma}

For the proof we refer to the appendix \todo. Next we distinguish into the two cases 

\textbf{Case 1} If the following holds 

\[
	\exists \mathcal O_1 \text{ satisfying unique provability}, x, sk, \vec{z_0} \neq \vec{z_0}'
\]

\todo{unsure if unique provability sufficient}

Note that for each tuple $(x,sk)$ every execution of the eval algorithm of the combiner 
corresponds to a specific DAG where the nodes correspond to an initial application of an algorithm 
$A_v$ on the inputs specified by the inwards-pointing edges to the node which produces a value $a_v$
and an algorithm $S_v$ that produces a value $s_v$ from $sk$ followed by a call to the an eval oracle. 
We can assume that since there is no need to make an oracle call with the same secret key and 
input value $a_v$ that all oracle calls are on fresh inputs. With $y[\mathcal O_0, \mathcal O_1]$
we denote the output of the combiner. Informally we want to argue that since every oracle call to $\mathcal O_0$
is on a fresh input we can just hard-code the vector to produce two different outputs. Then 
if we assume that the verification algorithm is canonical, i.e. gets all intermediate results 
via the proof, recomputes the results to ensure that the intermediate results are correct and then 
makes one verification query for every evaluation query, then both vectors must lead to accepting 
verifications by definition. Therefore the construction cannot be uniquely provable.

\textbf{Case 2} Since case $1$ doesnt hold 

\[
	\forall \mathcal O_1 \text{ satisfying unique provability}, x, sk, \vec{z_0} \neq \vec{z_0}'
\]

Then we show that $\BBComb{VRF}[\mathcal O_{x \mapsto 0}, \mathcal O_{\Prim{VRF}}]$. Since 
$\mathcal O_{x \mapsto 0}$ is uniquely provable this does meet the conditions and since the 
output doesn't depend on the answer of oracle $\mathcal O_0 = \mathcal O_{\Prim{VRF}}$
the adversary can simulate the execution of the combiner on a randomly sampled input $x$
by replying with random values until the combiner terminates. To do this however we must show that 
besides the oracle calls to eval no algorithm may depend on the secret key $sk$ (which the adversary 
doesn't know and therefore cannot simulate). For that note that if we interpret the computation 
as a DAG $G$ with the nodes being an oracle call to a candidates eval function with 
source corresponding to $x$ and the sink corresponding to the single application of an algorithm 
$A_y$ to the output of the predecessor nodes (corresponding to oracle calls) then we only 
need to show that $A_y$ can be computed efficiently by the adversary without access to the secret key. 
First assume that the algorithm depends on a source of randomness then the algorithm generates two different outputs
for the same inputs, by correctness, both of which need to be provable. This immediately leads to a 
break in the unique provability. Then assume, by contradiction that the output depends on a secret key, i.e. if 

\[
	\exists sk_1, sk_2, \vec z_0, x: y_1 = A_y(sk_1, x, \vec z_0, \mathcal O_1(\vec z_0)) \neq A_y(sk_2, x,\vec z_0, \mathcal O_1(\vec z_0)) = y_2
\]

We want to modify the insecure candidate scheme such that 

% \[
% 	\exists \lambda \; \Prob*{
% 		\begin{aligned}
% 			&pk_1 = pk_2\\ 
% 			\land \; & sk_1 = \widetilde{sk_1}\\ 
% 			\land \; & sk_2 = \widetilde{sk_2}
% 		\end{aligned}
% 		:
% 		\begin{aligned}
% 			(sk_1, pk_1) \sample \alg{Gen}{\BBComb{VRF}[\mathcal O_1, \mathcal O_2]}(1^\lambda)\\ 
% 			(sk_2, pk_2) \sample \alg{Gen}{\BBComb{VRF}[\mathcal O_1, \mathcal O_2]}(1^\lambda)
% 		\end{aligned}
% 	} > 0
% \]

We introduce the following notation, for any secret keys $sk_1, sk_2$ we have 

\[
		sk_1 \sqcap sk_2 \iff \exists x,\vec z_0, \vec z_1: A_y(sk_1,x,\vec z_0, \vec z_1) \neq A_y(sk_1,x,\vec z_0, \vec z_1)
\]

\[
		sk_1 \sqcup sk_2 \iff \exists pk: \Prob*{(sk_1,pk) \sample Gen} > 0 \land \Prob{(sk_2,pk) \sample Gen} > 0 \todo
\]

(We can remap to make one instance secure always.)

In that case we require by correctness that both $y_1$ and $y_2$ verify under $pk$
which breaks unique provability.
\newpage 

A necessary condition for a $(2;1)$-robust black-box $\VRF$-combiner is that 
for all candidate schemes $\vrf_0, \vrf_1$ that implement the $\VRF$ primitive there exist 
adversaries $\mathcal A_{\vrf_0}, \mathcal A_{\vrf_1}$ that break the security of $\vrf_1, \vrf_2$
respectively if there exists an adversary $\mathcal A_{\Pi^{\vrf_0, \vrf_1}}$ against 
the combiner instanciated with $\vrf_0, \vrf_1$. Formally 

\begin{equation*}
	\forall \vrf_1, \vrf_2 \; \Paren*{\exists \mathcal A_{\Comb\VRF{\vrf_0}{\vrf_1}} \text{ against }\Comb\VRF{\vrf_0}{\vrf_1} \implies \exists \mathcal A_{\vrf_0} \text{ against }\vrf_0 \land \exists \mathcal A_{\vrf_1} \text{ against }\vrf_1}
\end{equation*}

To show that no combiner exists we need to show that the negation of this statement holds. I.e. that there exist 
candidate schemes $\vrf_0, \vrf_1$ that implement a $\VRF$ at least one of which is secure such that there 
exists an adversary that breaks security of $\Comb\VRF{\vrf_0}{\vrf_1}$. Formally

\begin{equation*}
	\exists \vrf_0, \vrf_1 \; \Paren*{\exists \mathcal A_{\Comb\VRF{\vrf_0}{\vrf_1}} \text{ against }\Comb\VRF{\vrf_0}{\vrf_1} \land \Paren*{\vrf_0 \text{ secure }\lor \vrf_1 \text{ secure}}} 
\end{equation*}

Let $\Comb\VRF$ be an arbitrary PPT oracle machine, we distinguish into two cases based on 
whether the combiner actually uses both candidates to form the output during evaluation 
or otherwise whether one candidate is essentially ignored. 

\textbf{Case 1}: If the following holds 

% \begin{equation*}
% 	\exists \vrf_1, sk, x, \vec{z_0} \neq \vec{z_0}': y = g(sk, x, \vec{z_0}, \alg{Eval}{\vrf_2}(sk,x,\vec{z_0})) \neq g(sk, x, \vec{z_0}', \alg{Eval}{\vrf_2}(sk, x, \vec{z_0}')) = y'
% \end{equation*}

Then this means that at least one some secret key $sk \in \SK$ and input $x \in \bin^\lambda$ there exist values in place of evaluation of the first candidate scheme's outputs 
that produce a change in the evaluation algorithm of the combiner. Note that this alone doesn't give rise to two provable images under $x$ and $sk$ since $\vrf_1$ doesn't necessarily
need to be secure. Given a secure instance of a $\VRF$ we can however construct a new secure instance with the same outputs than $\vrf_1$ on the secret key $s_v(sk)$ for all $v \in V$ with $l(v) = 1$ that is still a secure VRF, 
see the construction \cref{fig:construction1}. We claim that then $\Comb\VRF{\vrf^0}{\widetilde{\vrf_1}}$ will both verify $(x,y,\pi), (x,y,\pi')$. This is because by the assumption above $y = z_t$ and $y' = z_t'$
and by definition of $\vrf^0$ every message there verifies, and by the construction of $\vrf_1$ every message there verifies and finally the relations $z_v = f_v(w_v)$ must also hold by construction.

% \begin{figure}[H]
% 	\centering 
% 	\begin{pchstack}[boxed]
% 		\procedure[linenumbering]
% 		{$\alg{Gen}{\widetilde{\vrf_1}}(1^\lambda)$}
% 		{
% 			(sk, pk) \sample \alg{Gen}{\vrf_1}(1^\lambda)\\ 
% 			\pcreturn (sk, pk)
% 		}
% 		\procedure[linenumbering]
% 		{$\alg{Eval}{\widetilde{\vrf_1}}(sk, x)$}
% 		{
% 			\pcif sk \in \{s_v(\widetilde{sk}) \mid v \in V, l(v) = 1\} \pcthen\\ 
% 			\t	(y, \pi) \sample \alg{Eval}{\vrf_1}(sk, x) \\
% 			\pcelse \\ 
% 			\t (y,\pi) \sample \alg{Eval}{\vrf_s}(sk, x)\\ 
% 			\pcreturn (y,\pi) 
% 		}
% 	\end{pchstack}
% 	\caption{Construction of the new secure instance of a $\VRF$}
% 	\label{fig:construction1}
% \end{figure}

\textbf{Case 2}: If the negation of the assumption in \textbf{Case 1} holds, i.e. 

\begin{equation*}
	\forall \vrf_1, sk, x, \vrf{z_0} \neq \vrf{z_0}': y = g(sk, x, \vec{z_0}, \alg{Eval}{\vrf_2}(sk,x,\vec{z_0})) = g(sk, x, \vec{z_0}', \alg{Eval}{\vrf_2}(sk, x, \vec{z_0}')) = y'
\end{equation*}

In that case the first candidate doesn't contribute to the final output for any secret key $sk$, input $x$. Formally there exists no sequence of possible outputs that would prompt a change in the output 
of the evaluation of the combiner. This does however mean that if the first candidate is the secure implementation that the combiner only depends on the insecure scheme. We show how this combiner is not 
pseudo-random. Formally we provide an adversary $\mathcal A_{\Comb\VRF{\vrf_0}{\vrf^0}}$ with non-negligible advantage in the PRF security experiment against $\Comb\VRF{\vrf_0}{\vrf^0}$. \todo{here we must show that none of the functions 
$f_v, v \in V \setminus \{s\}$ can depend on $sk$ if we exclude calls to oracles in them.} Then consider the following adversary $\mathcal A = (\mathcal A_1, \mathcal A_2)$ 

\begin{figure}[H]
	\centering
	\begin{pchstack}[boxed]
		\procedure[linenumbering]
		{Adversary $\mathcal A_1^{\mathcal O}(pk)$}
		{
			x \sample \bin^\lambda\\ 
			\sigma \gets x\\ 
			\pcreturn (x, \sigma)
		}
		\procedure[linenumbering]
		{Adversary $\mathcal A_2^{\mathcal O}(pk, \sigma, y)$}
		{
			x \gets \sigma\\ 
			\tilde y \gets g(x,\vec 0, \vec 0)\\ 
			\pcreturn y = \tilde y
		}
	\end{pchstack}
\end{figure}

\newpage

\newpage

\end{comment}

\section{Discussion and Extensions}

There are different ways to combine VRFs, for both unique provability and key-homomorphism 
they cannot be combined in a black-box fashion if we want to preserve properties such as 
unpredictability and therefore also randomness. On the other hand we manage to construct a 
robust black-box combiner that combines unbiasabile VRFs under guaranteed unique provability 
for all candidates. We find that the black-box combination of information theoretic properties 
generally seems hard if other properties such as unpredictability need to be preserved. 
While the class of all robust black-box combiners initially seems to be very broad we show that 
we can capture all polynomially bounded combiners using the notion of a robust black-box DAG combiner,
which allows a more systematic analysis of possible constructions by abstracting a graph structure 
in the computation executed by the implemented algorithms. We have not considered 
possible combiner constructions for stronger relaxations, such as computational unique provability 
or correctness with error, this is mainly because the resulting construction only offer weak 
guarantees that are insufficient for practical applications.

